%%% Hlavní soubor. Zde se definují základní parametry a odkazuje se na ostatní části. %%%

%% Verze pro jednostranný tisk:
% Okraje: levý 40mm, pravý 25mm, horní a dolní 25mm
% (ale pozor, LaTeX si sám přidává 1in)
\documentclass[12pt,a4paper]{report}
\setlength\textwidth{145mm}
\setlength\textheight{247mm}
\setlength\oddsidemargin{15mm}
\setlength\evensidemargin{15mm}
\setlength\topmargin{0mm}
\setlength\headsep{0mm}
\setlength\headheight{0mm}
% \openright zařídí, aby následující text začínal na pravé straně knihy
\let\openright=\clearpage

\let\sectionnobreak = \section
\makeatletter
\renewcommand{\section}{\@ifstar
                     \sectionStar%
                     \sectionNoStar%
}
\makeatother
\newcommand{\sectionNoStar}[1]{
	\pagebreak
	\sectionnobreak{#1}
}
\newcommand{\sectionStar}[1]{
	\pagebreak
	\sectionnobreak*{#1}
	\addcontentsline{toc}{section}{#1}
}

\makeatletter
\renewcommand{\paragraph}{%
  \@startsection{paragraph}{4}%
  {\z@}{1.5ex \@plus 1ex \@minus .2ex}{-1em}%
  {\normalfont\normalsize\bfseries}%
}
\makeatother

\usepackage[utf8]{inputenc}

%% Ostatní balíčky
\usepackage{graphicx}
\usepackage{amsthm}

\usepackage{upgreek}

\usepackage{listings}
\usepackage{tikz}
\usetikzlibrary{arrows}
\usepackage{xyling}
\usepackage[hyphens]{url}
\usepackage{microtype}
\usepackage{lmodern}
\usepackage{epstopdf}
\usepackage[printonlyused,nohyperlinks]{acronym}
\usepackage{paralist}
\usepackage{adjustbox}

\usepackage[center]{caption}

\usepackage{needspace}

\lstdefinelanguage{CSharp}[Sharp]{C}
{morekeywords={from,where,join,on,equals,orderby,descending,group,by,let,select},columns=flexible,showstringspaces=false}

\lstnewenvironment{code}[1][]%
{
   \noindent
   \minipage{\linewidth} 
   \vspace{0.5\baselineskip}
   \lstset{#1}}
{\endminipage}

\newcommand{\lstBreak}{\discretionary{\texttt{-}}{}{}}

\lstset{language=CSharp,basicstyle=\ttfamily}

%% Balíček hyperref, kterým jdou vyrábět klikací odkazy v PDF,
%% ale hlavně ho používáme k uložení metadat do PDF (včetně obsahu).
%% POZOR, nezapomeňte vyplnit jméno práce a autora.
\usepackage[unicode]{hyperref}   % Musí být za všemi ostatními balíčky
\hypersetup{pdftitle=.NET library for the MediaWiki API}
\hypersetup{pdfauthor=Petr Onderka}

\usepackage[all]{hypcap}

%%% Drobné úpravy stylu

% Tato makra přesvědčují mírně ošklivým trikem LaTeX, aby hlavičky kapitol
% sázel příčetněji a nevynechával nad nimi spoustu místa. Směle ignorujte.
\makeatletter
\def\@makechapterhead#1{
  {\parindent \z@ \raggedright \normalfont
   \Huge\bfseries \thechapter. #1
   \par\nobreak
   \vskip 20\p@
}}
\def\@makeschapterhead#1{
  {\parindent \z@ \raggedright \normalfont
   \Huge\bfseries #1
   \par\nobreak
   \vskip 20\p@
}}
\makeatother

% Toto makro definuje kapitolu, která není očíslovaná, ale je uvedena v obsahu.
\def\chapwithtoc#1{
\chapter*{#1}
\addcontentsline{toc}{chapter}{#1}
}

\def\secwithtoc#1{
\sectionnobreak*{#1}
\addcontentsline{toc}{section}{#1}
}

\newcommand\ApiParameterNoAmp[2]{\texttt{#1}~\texttt{=}~\texttt{#2}}
\newcommand\ApiParameter[2]{\ApiParameterNoAmp{#1}{#2}~\texttt{\&}}

\begin{document}

\hypersetup{pageanchor=false}

% Trochu volnější nastavení dělení slov, než je default.
\lefthyphenmin=2
\righthyphenmin=2

%%% Titulní strana práce

\pagestyle{empty}
\begin{center}

\large

Charles University in Prague

\medskip

Faculty of Mathematics and Physics

\vfill

{\bf\Large BACHELOR THESIS}

\vfill

\centerline{\mbox{\includegraphics[width=60mm]{img/logo}}}

\vfill
\vspace{5mm}

{\LARGE Petr Onderka}

\vspace{15mm}

% Název práce přesně podle zadání
{\LARGE\bfseries .NET library for the MediaWiki API}

\vfill

% Název katedry nebo ústavu, kde byla práce oficiálně zadána
% (dle Organizační struktury MFF UK)
Department of Theoretical Computer Science \linebreak and Mathematical Logic

\vfill

\begin{tabular}{rl}

Supervisor of the bachelor thesis: & Tomáš Petříček \\
\noalign{\vspace{2mm}}
Study programme: & Computer Science \\
\noalign{\vspace{2mm}}
Specialization: & General Computer Science \\
\end{tabular}

\vfill

% Zde doplňte rok
Prague 2012

\end{center}

\newpage

%%% Následuje vevázaný list -- kopie podepsaného "Zadání bakalářské práce".
%%% Toto zadání NENÍ součástí elektronické verze práce, nescanovat.

%%% Na tomto místě mohou být napsána případná poděkování (vedoucímu práce,
%%% konzultantovi, tomu, kdo zapůjčil software, literaturu apod.)

\openright

\noindent
I would like to thank to my supervisor, Tomáš Petříček,
for his help with writing this thesis.
I would also like to thank to my family for their unending support and patience during my studies.

\newpage

%%% Strana s čestným prohlášením k bakalářské práci

\vglue 0pt plus 1fill

\noindent
I declare that I carried out this bachelor thesis independently, and only with the cited
sources, literature and other professional sources.

\medskip\noindent
I understand that my work relates to the rights and obligations under the Act No.
121/2000 Coll., the Copyright Act, as amended, in particular the fact that the Charles
University in Prague has the right to conclude a license agreement on the use of this
work as a school work pursuant to Section 60 paragraph 1 of the Copyright Act.

\vspace{10mm}

\hbox{\hbox to 0.5\hsize{%
In Prague, date ............
\hss}\hbox to 0.5\hsize{%
signature of the author
\hss}}

\vspace{20mm}
\newpage

%%% Povinná informační strana bakalářské práce

\vbox to 0.5\vsize{
\setlength\parindent{0mm}
\setlength\parskip{5mm}

Název práce:
.NET knihovna pro MediaWiki API
% přesně dle zadání

Autor:
Petr Onderka

Katedra:  % Případně Ústav:
Katedra teoretické informatiky a matematické logiky 
% dle Organizační struktury MFF UK

Vedoucí bakalářské práce:
Mgr. Tomáš Petříček, University of Cambridge
% dle Organizační struktury MFF UK, případně plný název pracoviště mimo MFF UK

Abstrakt:
% abstrakt v rozsahu 80-200 slov; nejedná se však o opis zadání bakalářské práce

Wiki běžící na systému MediaWiki poskytují svým uživatelům API,
které lze použít k přístupu k dané wiki z počítačového programu.
Toto API je rozsáhlé, často se mění a může se lišit wiki od wiki,
takže může být náročné napsat knihovnu pro přístup k tomuto API.

Tato práce popisuje LinqToWiki,
knihovnu pro přístup k MediaWiki API ze C\# nebo jiných jazyků na platformě .NET.
Díky použití LINQu a generovaní kódu pomocí Roslynu,
kód napsaný s použitím této knihovny je čitelný, objevitelný, silně typovaný a flexibilní.

Klíčová slova: Wiki, C\#, LINQ, Generování kódu, Roslyn
% 3 až 5 klíčových slov

\vss}\nobreak\vbox to 0.49\vsize{
\setlength\parindent{0mm}
\setlength\parskip{5mm}

Title:
.NET library for the MediaWiki API
% přesný překlad názvu práce v angličtině

Author:
Petr Onderka

Department:
Department of Theoretical Computer Science and Mathematical Logic
% dle Organizační struktury MFF UK v angličtině

Supervisor:
Mgr. Tomáš Petříček, University of Cambridge
% dle Organizační struktury MFF UK, případně plný název pracoviště
% mimo MFF UK v angličtině

Abstract:
% abstrakt v rozsahu 80-200 slov v angličtině; nejedná se však o překlad
% zadání bakalářské práce

MediaWiki wikis provide their users an API, that can used to programmatically access the wiki.
This API is large, changes frequently and can be different from wiki to wiki,
so it can be a challenge to write a library for accessing the API.

This thesis describes LinqToWiki, a library
that can be used to access the MediaWiki API from C\# or other .NET languages.
Thanks to the use of LINQ and code generation through Roslyn,
code written using this library is readable, discoverable, strongly-typed and flexible.

Keywords: Wiki, C\#, LINQ, Code generation, Roslyn
% 3 až 5 klíčových slov v angličtině

\vss}

\newpage

%%% Strana s automaticky generovaným obsahem bakalářské práce. U matematických
%%% prací je přípustné, aby seznam tabulek a zkratek, existují-li, byl umístěn
%%% na začátku práce, místo na jejím konci.

\hypersetup{pageanchor=true}
\openright
\pagestyle{plain}
\setcounter{page}{1}
\tableofcontents

%%% Jednotlivé kapitoly práce jsou pro přehlednost uloženy v samostatných souborech
\chapwithtoc{Introduction}

Wiki websites running on the MediaWiki software (such as Wikipedia) offer an \ac{API}
for programmatic access to their database.
Since MediaWiki contains many functions, the \acs{API} is extensive too:
the core installation contains over seventy “modules” and more are available through extensions.
Each module represents a function available from the \acs{API}.
Modules accept parameters in the form of key-value pairs
and return a structured response in one of the supported formats, including \acs{XML}.

Because of the size of the API, accessing it from programming languages is not easy.
Two basic approaches are possible: static and dynamic:

\begin{itemize}
\item The dynamic approach is to create a thin library around the API modules:
let the user specify the names of the parameters and their types and return the response
in a dynamic manner, possibly as an associative array, or something like \acs{XML} \ac{DOM}.

This way, the user is responsible for the correctness of his query and for correct processing
of the response.
Also, it is hard to discover what the possible parameters are, what values can they have
and what form of response to expect.
This can lead to excessive use of the documentation or a “trial and error” approach.

\item The static approach is to create an extensive (or “thick”) library that contains methods tailored for every module,
each returning a different statically-typed result.

This way, many of the errors the user could make will result in a compile-time error
and the development environment can also advise the user what options are available.

But this approach is also inflexible: if the user wants to use something the library
was not made for, he cannot.
Differences like this can be caused by different versions of the software,
different sets of installed extensions, or just by different configuration.

Another question with the static approach is how to represent the parameters in code.
Most modules have many optional parameters, and so presenting them to the user
in an understandable manner might be a challenge.

One more problematic part is how to represent choosing which properties to return in the result.
A list of strings representing the chosen properties might be suitable for the dynamic approach, but not so much for the static one.
\end{itemize}

\medskip

This work introduces the LinqToWiki library and related tools that try to solve all those problems
using the C\# language and the .NET platform.

The dynamic vs. static issues are solved by automatically generating statically typed
code based on the metadata the API provides about itself.
The code generation is performed using Roslyn,
which is a new implementation of a compiler for the C\# language written in C\#.

The problems specific to the static approach are solved by using \ac{LINQ}:
a set of features of the C\# language and the .NET Framework,
that is useful for representing queries and their translation into another form.

The library also abstracts away some other aspects of the \ac{API}, like paging of the results.

\secwithtoc{Key contributions}

The key characteristics that make the LinqToWiki library novel, when compared with similar querying libraries are:

\begin{itemize}
\item Using different types in different LINQ operators
(specifically in \lstinline{select}, \lstinline{where} and \lstinline{orderby})
\item Allowing different operators for different queries
(i.e. allowing sorting and generator queries only for some modules and disallowing sorting twice in one query)
\item Using code generation to achieve statically-typed queries
\end{itemize}

Another important aspect of this work was creating a patch for MediaWiki to describe results of \ac{API} modules,
which was accepted by MediaWiki developers and is now running on all Wikimedia wikis, including Wikipedia.

\secwithtoc{Structure of this work}
Chapter~\ref{goal} explains what is the goal of the LinqToWiki library.
Chapter~\ref{background} describes libraries and \acp{API}
that were important for creating this work.
Specifically, those are the MediaWiki API (Section~\ref{mediawiki}),
LINQ (Section~\ref{linq}) and Roslyn (Section~\ref{roslyn}).
Chapter~\ref{mw improvements} talks about changes that were made to MediaWiki,
to make LinqToWiki possible.
Chapter~\ref{ltw} describes the library itself, and the projects it consists of.
Chapter~\ref{future work} mentions some ways in which the library could be further improved.
Chapter~\ref{related} describes other libraries for accessing MediaWiki API.

\medskip

Appendix~\ref{cd} describes the contents of the enclosed CD.

\chapter{Problem analysis}
\label{goal}

The goal of the LinqToWiki library is to be able to express requests using the MediaWiki \ac{API}
in a way that is readable, discoverable, checked by the compiler for correctness as much as possible and also flexible with regards to changes.

This is achieved by generating classes specific for each module and using them in LINQ queries.

\paragraph{Querying data in C\#}

\ac{LINQ} is a way of querying various data sources from the C\# language.
The two most commonly used variants are LINQ to Objects, and various versions of LINQ for \ac{SQL} databases.
LINQ to Objects is used for querying in-memory data, like arrays.
There are several widely-used libraries for accessing \ac{SQL} databases using LINQ, including LINQ to SQL, LINQ to Entities and NHibernate.

In all versions of \ac{LINQ}, the queries look the same. For example:

\begin{lstlisting}
from product in products
where product.Price > 500
   && product.InStock
join category in categories on product.Category equals category
orderby product.Price
select product.Name
\end{lstlisting}

A query like this is translated into a sequence of method calls that take their parameters in the form of lambda expressions.%
\footnote{
Lambda expressions are a way to create unnamed (anonymous) functions.
Parameter list of the function is specified on the left side of the \lstinline{=>} operator,
result of the function is computed by the expression on its right side.
Types of the parameters and the return type do not have to be specified, if they can be inferred.
}
For example, the \lstinline{where} part of the above query is translated into:

\begin{lstlisting}
products.Where(product => product.Price > 500 && product.InStock)
\end{lstlisting}

The commonalities between LINQ to Objects and SQL LINQ libraries are that the full range of operators is available
and that all properties of the queried type are available in all of them.

\paragraph{Querying MediaWiki}

The situation with the MediaWiki \ac{API} is different in several ways:

\begin{enumerate}
\item It does not support queries represented by many of the LINQ operators, including \lstinline{join} and \lstinline{group by}.
\item Some of the modules do not support sorting, some do. Of those that do support sorting, some allow specifying the sort key, others only the direction.
\item The sets of properties that are available for filtering, sorting and selecting are all different.

For example, the \texttt{blocks} module does support specifying the direction, but not the property to sort by.
Its set of properties that can be used for filtering includes \lstinline{Ip}
(the \ac{IP} address of the blocked user), but not \lstinline{ById}.
On the other hand, properties that can be used for selecting include \lstinline{ById}
(ID of the user who performed the block), but not \lstinline{Ip}.

\item There are modules used for queries about a set of pages. Those pages can be from a hard-coded list or a result from some other module.

An example of such module is \texttt{categories}, which returns the list of categories of each page
specified in the input list.

\item There are also parameters that do not fit into the LINQ model well. Some of them are required, some are not.

These parameters usually modify which items appear in the output,
but they are not just a filter from some larger lists.
Examples of such parameters are the required parameter \texttt{search} of the \texttt{search} module
(which specifies what string to search for)
or the optional parameter \texttt{title} of the \texttt{categorymembers} module
(which decides what category to enumerate;
this parameter is not required, because the parameter \texttt{pageid} can be used instead of it).
\end{enumerate}

The goal is to be able to represent all valid queries, while invalid queries should cause a compile-time error.

Specifically, unsupported operators (like \lstinline{join} and \lstinline{group by}) should cause an error for all modules,
while the \lstinline{orderby} clause should cause an error only for the modules that do not support sorting.

Also, all operators should support only those properties that are actually supported by the \ac{API}.
So, for example for the \texttt{blocks} module, the following query should compile and execute fine:

\begin{lstlisting}
from block in wiki.Blocks()
where block.Ip == "8.8.8.8"
orderby block descending
select block.ById
\end{lstlisting}

This is because
\begin{compactitem}
\item limiting the query by the blocked \ac{IP} address,
\item sorting without specifying the key and
\item selecting the ID of the user who performed the block
\end{compactitem}
are all allowed, while the following query should cause three errors:

\begin{lstlisting}
from block in wiki.Blocks()
where block.ById == 1234
orderby block.Expiry descending
select block.Ip
\end{lstlisting}

Here,
\begin{compactitem}
\item limiting by the ID of the user who performed the block,
\item sorting by the expiration date and
\item selecting the \ac{IP} address
\end{compactitem}
are all impossible.
(Actually selecting the \ac{IP} address of the blocked user is possible,
but the information is contained in properties with different names.)

\section{Alternatives}
\label{alternatives}

Using statically typed methods and custom LINQ provider is not the only way a library like this
could be built in C\#.
Some of the alternatives include:

\begin{itemize}
\item Using strings for everything

This is probably the simplest way to write a library for MediaWiki API.
It means parameters and their names (including those specifying which module to use) are specified
as a collection of string key-value pairs.

For example, the query from the previous section would look something like:

\begin{lstlisting}
wiki.Query(new Dictionary<string, string>
    {
        { "list", "blocks" },
        { "bkip", "8.8.8.8" },
        { "bkdir", "older" },
        { "bkprop", "byid" }
    })
\end{lstlisting}

The main disadvantage of this approach is there is no checking of the query:
any mistake in the query won't be found out until the query is actually executed.
It also means autocompletion cannot help the user in finding out which parameters are available.

In such implementation, the result of the query will be also some sort of dictionary
(indexed by names of result properties), which has the same kind of problems as using strings for parameter names.

Some modifications are possible, for example having a special parameter for the module prefix
(\texttt{bk} in the above example), which would make the code less repetitive.
But such changes will not fix the fundamental problems with this approach.

\item Using \lstinline{dynamic}

C\# 4.0 supports dynamic typing using the special type \lstinline{dynamic}.
When using that, C\# acts similarly as dynamic languages.
This means methods and their parameters do not have to be declared before they can be used.

For example, the same query used previously could look like this using \lstinline{dynamic}:

\begin{lstlisting}
dynamic wiki = ...;
wiki.ListBlocks("bk", ip:"8.8.8.8", dir:"older", prop:"byid")
\end{lstlisting}

For this query to work,
the library does not need any knowledge about the \texttt{blocks} module or its parameters.

This query is much more succinct than the string version
(although part of that is because of extracting the prefix,
which could be done with strings too, as mentioned).
But it suffers from the same issues: mistakes cannot be detected at compile time
and autocomplete will not be able to help.

In this case, the returned object will be also dynamic,
with the same advantages and disadvantages.

One interesting consequence of using \lstinline{dynamic} is that it means LINQ cannot be used.
Some of the reasons for this are that dynamically invoked methods cannot have lambdas as their arguments
(at least not directly) and that expression trees do not support \lstinline{dynamic}.

\item Manually written code for each module

Here, code for each module would be written manually, instead of using code generation.
This would be a significant amount of code which would have to be updated with every change of the API.
Another problem is that each wiki can have different set of extensions, which add their own API modules,
so distinct versions for some wikis would have to be maintained.

The biggest advantage of this approach is that a human can understand the documentation of a module,
so they can write code for a module better than a code generator such as LinqToWiki.Codegen could.
This is because not all details of how a module works can be encoded into its machine readable description.

\item Using \lstinline{IQueryable} instead of custom LINQ provider

The \lstinline{IQueryable} interface is often used for implementing custom \ac{LINQ} sources,
especially those for querying \ac{SQL} databases.

The biggest disadvantage of using this approach is that \lstinline{IQueryable}
supports all LINQ operators and all operators use the same set of properties.
This is not suitable for MediaWiki \ac{API},
because it does not support several LINQ operators,
because ordering is available only for some modules
and because filtering, sorting and projection all use different sets of properties.

\end{itemize}

\chapter{Background}
\label{background}
\sectionnobreak{MediaWiki API}
\label{mediawiki}

MediaWiki \cite{mediawiki} is an open source wiki system.
It is written in the \acs{PHP} programming language and uses a relational database to store its data,
usually MySQL.
It is maintained by the Wikimedia Foundation,
who also runs some of the biggest wiki sites (or “wikis” for short),
including Wikipedia and Wiktionary.
It is also used by many others, including Wikia, who runs many small wikis for various interests and
the unofficial wiki of the Faculty of Mathematics and Physics, $\upomega\upiota\upkappa\upiota$.matfyz.cz.

\paragraph{Acessing MediaWiki}

There are several ways to programmatically access the database of some MediaWiki wiki.
First, it is possible to directly access the database using \ac{SQL}.
This usually requires access to the server that runs the database, so it is not available in many cases.
For Wikimedia wikis, this data is accessible through the Toolserver.
The Toolserver is a cluster of servers that contains replicated read-only databases of Wikimedia wikis,
primarily for use by various tools.
These copies of the databases do not contain article texts.
The Toolserver is run by Wikimedia Deutschland.

Mostly specific to Wikimedia wikis is also another option: database dumps. These are files that contain
dumps of some tables of the wikis. Their disadvantages are that the newest dump is usually several days or weeks old
and that the files can be huge, which is impractical for getting information about a small number of pages.

Last, but not least, is the MediaWiki API \cite{mediawiki-api}.
It can be used to remotely access any MediaWiki wiki (unless the \ac{API} is disabled in the configuration)
using \ac{HTTP}.

\paragraph{API request and response}

Parameters for an \ac{API} request are given in the query string of a GET request or in the body of a POST request
(modules that perform modifications require the use of POST).
The body of the POST request is usually formatted as \path{application/x-www-form-urlencoded},
but file uploads require the use of \path{multipart/form-data}.

Some parameters can accept multiple values at once.
In these cases, the values are separated by a pipe character (\texttt{|}).
For example, to retrieve information about the articles “LINQ” and “MediaWiki” from some module,
the parameter would be \texttt{titles=LINQ|MediaWiki}.

There are some parameters that are common to many modules:

\begin{itemize}
\item The \texttt{prop} parameter is used to determine what properties will be present in the response.
The values of this parameter don't map directly to the properties of the response,
so for example specifying \texttt{prop=ids} might cause the property \texttt{pageid} to appear in the result.
\item The \texttt{sort} parameter decides which property will be used to order the results of the query.
\item The \texttt{dir} parameter is used to determine the order of the results:
whether it should be ascending or descending.
\end{itemize}

The response can be in one of the several available formats, the most widely used ones are \ac{XML} and \ac{JSON}.

The representation of most data types in the response is standard:
\texttt{string}s are formatted as strings, \texttt{integer}s as decimal numbers,
\texttt{timestamp}s are formatted according to ISO 8601.
Only \texttt{boolean}s have an unexpected representation:
if the property is \texttt{false}, it is not present in the result at all,
and if its value is \texttt{true}, it is represented as an empty string.

If there is some problem executing a request, for example if a parameter has an invalid value,
a warning will be returned along with the result of the operation.
In the case of a fatal problem, such as when the user doesn't have the right to perform an action,
an error is returned, without any results.

\paragraph{API modules}

The \ac{API} is divided into modules and there are two kinds of modules:
“normal” modules (called “non-query modules” in this work) and query modules.

Non-query modules are usually used to perform some action.
For example the \texttt{edit} module can be used to edit a page
and the \texttt{block} module can be used to block another user (it can be used only by users with sufficient privileges).

Query modules are used for retrieving information about the wiki. There are three types of query modules:

\begin{itemize}
\item \texttt{list} modules: Return contents of various lists.
For example the \texttt{all\-categories} module can be used to list all categories on a wiki,
while the \texttt{categorymembers} module can be used to list members of a certain category.
\item \texttt{prop} modules: Return information about a set of pages.
For example, the \texttt{categories} module can be used to retrieve the categories for each page in a given set.
\item \texttt{meta} modules: Return meta information that are not directly associated with pages.
For example the \texttt{userinfo} module can be used to retrieve information about the currently logged-in user.
\end{itemize}

For \texttt{prop} modules, the set of pages they operate on can be specified directly using page titles or page IDs.

Another option is to use some other module (usually a \texttt{list} module) as a so called “generator”.
This way, one can for example retrieve all categories of pages in a specific category,
by using the \texttt{categorymembers} module as a generator for the \texttt{categories} module.

Because more than one module can be used in one request,
the parameters for each module are distinguished by using prefixes.
For example, the prefix for the \texttt{categorymembers} module is \texttt{cm}.
So, setting its \texttt{limit} parameter to the value of~5 can be achieved by
adding \texttt{cmlimit=5} to the query string of a GET request or to the body of a POST request.

The \ac{API} is also extensible: MediaWiki extensions can add their own modules and modify some behavior of existing modules.

Examples of \ac{API} requests and responses are
in Figures~\ref{API edit example} through \ref{API categories generator example}.

\begin{figure}[p]
\texttt{POST:}
\texttt{http://en.wikipedia.org/w/api.php}

\ApiParameter{format}{xml}
\ApiParameter{action}{edit}
\ApiParameter{token}{\%2B\%5C}
\ApiParameter{title}{Wikipedia:Sandbox}
\ApiParameter{section}{new}
\ApiParameter{sectiontitle}{Hello}
\ApiParameter{text}{Hello\%20world!\textasciitilde\textasciitilde\textasciitilde\textasciitilde}
\ApiParameterNoAmp{summary}{greeting\%20the\%20world}

\begin{lstlisting}[language=xml]
<?xml version="1.0"?>
<api>
  <edit result="Success" pageid="16283969"
         title="Wikipedia:Sandbox" contentmodel="wikitext"
         oldrevid="526507943" newrevid="526508868"
         newtimestamp="2012-12-05T13:10:24Z" />
</api>
\end{lstlisting}

\caption[API request to add a new section to a page and its response]{
\ac{API} request to add a “Hello” section to the page \texttt{Wikipedia:\allowbreak Sandbox}, \\
using the non-query module \texttt{edit} and its response}
\label{API edit example}
\end{figure}

\begin{figure}[p]
\texttt{GET:}
\texttt{http://en.wikipedia.org/w/api.php}~\texttt{?}
\ApiParameter{format}{xml}
\ApiParameter{action}{query}
\ApiParameter{list}{categorymembers}
\ApiParameter{cmtitle}{Category:Query\%20languages}
\ApiParameter{cmprop}{title}
\ApiParameter{cmtype}{page}
\ApiParameter{cmdir}{descending}
\ApiParameterNoAmp{cmlimit}{5}

\begin{lstlisting}[language=xml]
<?xml version="1.0"?>
<api>
  <query>
    <categorymembers>
      <cm ns="0" title="Yahoo! Query Language" />
      <cm ns="0" title="XQuery" />
      <cm ns="0" title="XPath" />
      <cm ns="0" title="XBase++" />
      <cm ns="0" title="XBase" />
    </categorymembers>
  </query>
  <query-continue>
    <categorymembers cmcontinue="page|57454220434f5645524147..." />
  </query-continue>
</api>
\end{lstlisting}

\caption[API request to list pages from a category and its response]{
\ac{API} request to list the last five pages \\
from the category \texttt{Category:\allowbreak Query languages} \\
using the \texttt{list} module \lstinline{categorymembers}
and its response}
\label{API categorymembers example}
\end{figure}

\begin{figure}[p]
\texttt{GET:}
\texttt{http://en.wikipedia.org/w/api.php}~\texttt{?}
\ApiParameter{format}{xml}
\ApiParameter{action}{query}
\ApiParameter{titles}{LINQ|MediaWiki}
\ApiParameter{prop}{categories}
\ApiParameterNoAmp{cllimit}{5}

\begin{lstlisting}[language=xml]
<?xml version="1.0"?>
<api>
  <query>
    <pages>
      <page pageid="11904240" ns="0" title="LINQ" />
      <page pageid="323710" ns="0" title="MediaWiki">
        <categories>
          <cl ns="14" title="Category:2002 software" />
          <cl ns="14" title="Category:All articles with spec..." />
          <cl ns="14" title="Category:All articles with unso..." />
          <cl ns="14" title="Category:Articles with specific..." />
          <cl ns="14" title="Category:Articles with unsource..." />
        </categories>
      </page>
    </pages>
  </query>
  <query-continue>
    <categories clcontinue="323710|Commons_category_template..." />
  </query-continue>
</api>
\end{lstlisting}

\caption[API request to list categories of given pages and its response]{
API request to list categories of the pages \texttt{LINQ} and \texttt{MediaWiki} \\
using the \texttt{prop} module \texttt{categories} and its response}
\label{API categories titles example}
\end{figure}

\begin{figure}[p]
\texttt{GET:}
\texttt{http://en.wikipedia.org/w/api.php}~\texttt{?}
\ApiParameter{format}{xml}
\ApiParameter{action}{query}
\ApiParameter{generator}{categorymembers}
\ApiParameter{gcmtitle}{Category:Query\%20languages}
\ApiParameter{gcmtype}{page}
\ApiParameter{gcmdir}{descending}
\ApiParameter{gcmlimit}{2}
\ApiParameter{prop}{categories}
\ApiParameterNoAmp{cllimit}{5}

\begin{lstlisting}[language=xml]
<?xml version="1.0"?>
<api>
  <query-continue>
    <categorymembers gcmcontinue="page|58504154480a585041544..." />
    <categories clcontinue="23742879|Query_languages" />
  </query-continue>
  <query>
    <pages>
      <page pageid="22672444" ns="0" title="Yahoo! Query Language">
        <categories>
          <cl ns="14" title="Category:Query languages" />
          <cl ns="14" title="Category:World Wide Web stubs" />
          <cl ns="14" title="Category:Yahoo! Development" />
        </categories>
      </page>
      <page pageid="23742879" ns="0" title="XQuery">
        <categories>
          <cl ns="14" title="Category:4GL" />
          <cl ns="14" title="Category:Functional languages" />
        </categories>
      </page>
    </pages>
  </query>
</api>
\end{lstlisting}

\caption[API request to list categories of pages from a category and its response]{
\ac{API} request to list categories of pages
from the end of \texttt{Category:\allowbreak Query languages}
using the \texttt{list} module \texttt{categorymembers} as a generator
for the \texttt{prop} module \texttt{categories}
and its response}
\label{API categories generator example}
\end{figure}

\subsection{Paging}
\label{mw paging}

Because the results of the \ac{API} queries can contain thousands and sometimes even millions of entries,
the responses are limited.
For most modules, the default limit (when it is not specified as a parameter) is ten entries per page
and the default maximum is 500~entries for normal users.
For users with the \texttt{apihighlimits} right, the limits are raised, usually to 5000~entries per page.

In the \texttt{limit} parameter, one can specify either the exact value,
or the special value \texttt{max}, which means the maximum allowed for the current user.

To get the data from the following page, one has to use a value specified in the \texttt{query-continue}
element in the result (see Figure \ref{API categorymembers example} again).
The value in this element is a transparent identifier of the next page.

The advantage of this system when compared with the conventional paging systems of numbering pages
or using numeric offsets is that it avoids missing entries and duplicates when the result
changes while retrieving the pages.

The \ac{API} has no notion of transactions, so it is not possible to get fully consistent results
that would correspond to an exact moment in time.
But thanks to this paging system, one can be certain that an entry that should be in the result set
during retrieving of all of the pages will actually be present in the result set exactly once.

\paragraph{Paging with \texttt{prop} modules}

The situation gets more complicated when using a \texttt{prop} module with another module as a generator
(see Figure~\ref{API categories generator example} again for an example).
That is because both modules have their own paging.

When such a request is made, the first response will contain a limited number of items from the generator
and a limited number of results from the \texttt{prop} module for those items.
To retrieve the next set of items from the generator, one has to use the \texttt{query-continue} for the generator
(called “primary paging” in this work).
To retrieve the next set of results for the items from the first result,
one has to use the \texttt{query-continue} for the \texttt{prop} module (called “secondary paging” here).

For an example, see Figure \ref{paging}.
It shows how the paging might work when using the \texttt{allpages} module as a generator,
together with the \texttt{prop} module \texttt{categories}.
The \texttt{query-continue} elements in the figure are just examples,
actual values would be more complicated.

\begin{figure}[tbp]
\begin{center}
\begin{tikzpicture}
\path (0,17) node(11) {
\begin{adjustbox}{minipage=160pt,fbox}
\begin{lstlisting}[language=xml]
<query>
  <pages>
    <page title="A" />
    <page title="B">
      <categories>
        <cl title="X" />
        <cl title="Y" />
      </categories>
    </page>
  </pages>
  <query-continue>
    <categorymembers
      gcmcontinue="A|X" />
    <allpages
      apcontinue="C" />
  </query-continue>
</query>
\end{lstlisting}
\end{adjustbox}
} (6.9,17) node(12) {
\begin{adjustbox}{minipage=160pt,fbox}
\begin{lstlisting}[language=xml]
<query>
  <pages>
    <page title="A">
      <categories>
        <cl title="X" />
        <cl title="Y" />
      </categories>
    </page>
    <page title="B" />
  </pages>
  <query-continue>
    <categorymembers
      gcmcontinue="A|Z" />
    <allpages
      apcontinue="C" />
  </query-continue>
</query>
\end{lstlisting}
\end{adjustbox}
} (11,17) node(13) {\dots};

\path (0,6) node(21) {
\begin{adjustbox}{minipage=160pt,fbox}
\begin{lstlisting}[language=xml]
<query>
  <pages>
    <page title="C">
      <categories>
        <cl title="X" />
        <cl title="Y" />
      </categories>
    </page>
    <page title="D" />
  </pages>
  <query-continue>
    <categorymembers
      gcmcontinue="C|Z" />
    <allpages
      apcontinue="E" />
  </query-continue>  
</query>
\end{lstlisting}
\end{adjustbox}
} (6.9,6) node(22) {
\begin{adjustbox}{minipage=160pt,fbox}
\begin{lstlisting}[language=xml]
<query>
  <pages>
    <page title="C">
      <categories>
        <cl title="Z" />
      </categories>
    </page>
    <page title="D">
      <categories>
        <cl title="X" />
      </categories>
    </page>
  </pages>
  <query-continue>
    <allpages
      apcontinue="E" />
  </query-continue>  
</query>
\end{lstlisting}
\end{adjustbox}
};

\path (0,0) node(31) {\dots};

\draw[-triangle 45] (11) -- (12);
\draw[-triangle 45] (12) -- (13);
\draw[-triangle 45] (11) -- (21);
\draw[-triangle 45] (21) -- (22);
\draw[-triangle 45] (21) -- (31);

\end{tikzpicture}
\end{center}

\caption{Example of primary and secondary paging}
\label{paging}

\end{figure}

\paragraph{Paging with the \texttt{revisions} module}

The situation is even more complicated with the \texttt{prop} module \texttt{revisions}.
It can be used to retrieve information about revisions of pages, including their text
and it is the only module that can be used to get the text of a set of pages.

For this module, not specifying the \texttt{limit} parameter means that only the most
recent revision will be shown and no \texttt{query-continue} will be present.
Also, when \texttt{limit} is specified, the module can operate only on one page at a time,
so for example one has to set the \texttt{limit} of a module used as a generator to~1.

But for other modules, when no \texttt{limit} parameter is specified, a default value is used (usually 10)
and a \texttt{query-continue} element is present in the response, to access the remaining items.

\subsection{The \texttt{paraminfo} module}
\label{paraminfo}

The \texttt{meta} query module \texttt{paraminfo} is of a special importance for this work.
It can be used to retrieve information about modules,
which is necessary for generating code to access those modules in a static fashion.

Prior to an improvement implemented in our work,
the \texttt{paraminfo} module provided only general information about the module
and, most importantly, for each parameter information about its data type and a short description.

The data type of a parameter is either a simple type (e.g. \texttt{integer} or \texttt{string}),
or an enumeration of possible values.
And the description is useful as a documentation for the generated code.

A shortened example of a response from the \texttt{paraminfo} module
for the \texttt{categorymembers} module is in Figure~\ref{paraminfo sample}.

For code generation in LinqToWiki, another piece of information is necessary:
knowing the properties of the response and how do they map to the values of the \texttt{prop} parameter
(which determines what properties are included in the response).
For information about how we added them, see Chapter~\ref{mw improvements}.

\begin{figure}[p]

\begin{lstlisting}[language=]
<module name="categorymembers" prefix="cm" querytype="list"
  generator="" listresult="" description="List all pages in a ...">
  <parameters>
    <param name="title" type="string" 
      description="Which category to enumerate (required). ..." />
    <param name="pageid" type="integer" 
      description="Page ID of the category to enumerate. ..." />
    <param name="prop" default="ids|title" multi=""
      description="What pieces of information to include ...">
      <type>
        <t>ids</t>
        <t>title</t>
        <t>sortkey</t>
        <t>sortkeyprefix</t>
        <t>type</t>
        <t>timestamp</t>
      </type>
    </param>
    <param name="namespace" multi="" type="namespace"
      description="Only include pages in these namespaces" />
    <param name="continue" type="string"
      description="For large categories, give the value ..." />
    <param name="limit" default="10" max="500" type="limit" 
      description="The maximum number of pages to return." />
    <param name="sort" default="sortkey"
      description="Property to sort by">
      <type>
        <t>sortkey</t>
        <t>timestamp</t>
      </type>
    </param>
    <param name="dir" default="ascending"
      description="In which direction to sort">
      <type>
        <t>ascending</t>
        <t>descending</t>
      </type>
    </param>
  </parameters>
</module>
\end{lstlisting}

\caption{Shortened response of the \texttt{paraminfo} module \\ for the \texttt{category\-members} module}
\label{paraminfo sample}
\end{figure}
\section{LINQ and expression trees}
\label{linq}

\acs{LINQ}, short for Language Integrated Query, is a feature of the C\# programming language%
\footnote{\ac{VB.NET} also supports \ac{LINQ}, with slightly different syntax
and capabilities, but uses the same types.}
and the .NET Framework that can be used for querying various data sources
and appeared in the version 3.0 of the language \cite{cs-in-depth}.
It uses higher-order functions and lambda expressions to achieve a readable declarative syntax.

\ac{LINQ} consists of a set of so called “standard query operators”:
methods that are used to perform the query operations on a given source.
Also, a special syntax (called “query expressions”), similar to \ac{SQL} queries, is available
for some of those operators.
The compiler translates a query expression into a set of calls to standard query operators,
using lambda expressions and anonymous types.

Anonymous types are types that do not have to be explicitly declared;
they are used in similar situations as tuples in functional programming.
An instance of an anonymous type is created by using the \lstinline{new} keyword
without specifying the type of the object to create.

\paragraph{Query translation}

For example the following query expression (as seen in Chapter \ref{goal}):

\nopagebreak

\begin{code}
from product in products
where product.Price > 500
   && product.InStock
join category in categories on product.Category equals category
orderby product.Price
select product.Name
\end{code}

Is translated into the following method calls:

\begin{code}
products
    .Where(product => product.Price > 500 && product.InStock)
    .Join(
        categories,
        product => product.Category,
        category => category,
        (product, category) => new { product, category })
    .OrderBy(t => t.product.Price)
    .Select(t => t.product.Name)
\end{code}

The parameter~\lstinline,t, is called “transparent identifier”. It is used to transfer a set of variables from one method call to another.

The \ac{LINQ} library also contains methods that do not have a corresponding representation in query expressions. Some examples of those are \lstinline,Aggregate(),, \lstinline,Sum(), and \lstinline,ToList(),.

\paragraph{Relationship with functional programming}

Many of the basic query operators also correspond to well-known higher-order functions from functional programming. See Figure \ref{LINQ methods} for comparison of some of the \ac{LINQ} query operators, query expression clauses, and higher-order functions.

\begin{figure}[htbp]
\begin{tabular}{lll}
query operator & query expression clause & functional name \\
\hline
\lstinline,Select(), & \lstinline,select,, \lstinline,let, & map \\
\lstinline,Where(), & \lstinline,where, & filter \\
\lstinline,SelectMany(), & second and following \lstinline,from, & bind \\
\lstinline,Aggregate(), & & fold \\
\lstinline,Join(), & \lstinline,join, & \\
\lstinline,OrderBy(),,\cr \lstinline,OrderByDescending(), & \lstinline,orderby, & sort \\
\lstinline,GroupBy(), & \lstinline,group by, & \\
\lstinline,Sum(), \\
\lstinline,First(), & & head \\
\lstinline,ToList(), \\
\end{tabular}

\caption{Comparison between \ac{LINQ} query operators, query expression clauses
and functions in functional languages}
\label{LINQ methods}
\end{figure}

\paragraph{Expression trees}

Usually, lambda expressions are compiled into normal methods and passed to the query operator methods as delegates
(which are similar to function pointers in~C or first-class functions in functional languages).
But this would not be suitable for querying of sources that are not in-memory collections.
This is because the query has to be translated into another form,
like an \ac{SQL} query or a set of parameters for the MediaWiki \ac{API}.

Because of this, a lambda expression in C\# can be also compiled into another form:
an expression tree.
Expression tree is an object that represents the given lambda expression
in a form similar to an abstract syntax tree.
This object can be programmatically accessed and manipulated,
which allows translation of \ac{LINQ} queries into other forms, such as \ac{SQL} queries.
An expression tree can also be compiled into a delegate and then executed.

For an example of an expression tree, see Figure \ref{Expression tree}.

\begin{figure}[htbp]
\begin{center}
\Treek[2]{4}{
 & & & \K{Binary: AndAlso} \B{dll}_{\textnormal{Left}} \B{dr}^{\textnormal{Right}} \\
 & \K{Binary: GreaterThan} \B{dl}_{\textnormal{Left}} \B{dr}^{\textnormal{Right}} & & & \K{Member: InStock} \B{d}^{\textnormal{Expression}} \\
 \K{Member: Price} \B{d}_{\textnormal{Expression}} & & \K{Constant: 500} & & \K{Parameter: product} \\
\K{Parameter: product}
}
\end{center}

\caption[Example body of expression tree]
{The body of the expression tree for the lambda expression \lstinline,product => product.Price > 500 && product.InStock,}
\label{Expression tree}
\end{figure}

\paragraph{Implementing the LINQ pattern}

The .NET Framework contains two implementations of the query expression pattern:
the interfaces \lstinline,IEnumerable<T>, and \lstinline,IQueryable<T>,.
This means that any object that implements one of these two interfaces can be used in a \ac{LINQ} query.

These two types implement the query expression pattern completely,
so they can be used with any \ac{LINQ} operator.
Other custom types can implement only part of the query expression pattern,
which would mean only a subset of the \ac{LINQ} operators are available for such types.

The \lstinline,IEnumerable<T>, interface usually represents an in-memory collection,
so its implementation of the \ac{LINQ} operators use delegates.
The \lstinline,IQueryable<T>, interface is usually used to represent a remote collection
(such as a table in a relational database),
so its version of the \ac{LINQ} operators use expression trees.

The \lstinline,IQueryable<T>, interface does not perform any translation of expression trees
into the target query language.
Instead, it combines the whole query into one expression tree,
which is then passed to an implementation of \lstinline,IQueryProvider,.

The query provider is then responsible for processing the expression tree
and translating it into its target query language.
If the query is not valid, the query provider will throw an exception at runtime.
\section{Roslyn}
\label{roslyn}

Microsoft Roslyn is a new implementation of the C\# compiler written in C\#
(and a \ac{VB.NET} compiler written in \ac{VB.NET}) \cite{roslyn}.
Its main distinguishing characteristic is that it is “open”:
it can be used for example to convert between text and a syntax tree,
to manipulate the syntax tree or to retrieve semantic information.

It also integrates itself into the Microsoft Visual Studio \ac{IDE},
where it can be used to perform
custom refactoring actions or to produce custom errors and warnings at compile-time.

Roslyn is currently under development and so far it had three public releases.
All of them were in the form of \ac{CTP},
the first one from October 2011, the second one from June 2012 and the third one from September 2012.

In the \acp{CTP}, the syntactic part of the library in completely implemented,
so for example the syntax tree can represent any construct of C\#
and any syntax tree can be translated to and from source code.
On the other hand, the semantic part of the library is not fully implemented,
which means that for example some syntax trees will not successfully compile,
even if they represent valid C\# code.

Because of its close relation with Visual Studio,
Roslyn syntax tree is able to represent every feature of C\# with down to character precision.
This includes “trivia”: parts of code that are not significant for the compiler,
such as whitespace and comments.

Trivia can also be “structured”, that is, it can form a small syntax tree of its own.
An example of structured trivia are \ac{XML} documentation comments,
that can be used to provide documentation for a piece of code,
which can then be automatically processed.

For an example of a Roslyn syntax tree, see Figure~\ref{Syntax tree}

\begin{figure}[htbp]

\begin{lstlisting}
public abstract CategoryInfoResult CategoryInfo { get; }
\end{lstlisting}

\begin{center}
\Treek[.5]{4}{
& & & & \K{Property}\Below{Declaration} \B[-5]{dllll} \B[-5]{dll} \B[-5]{d} \B[-5]{drr} \B[-5]{drrrr} \\
\K{Public}\Below{Keyword} & & \K{Abstract}\Below{Keyword} & & \K{Identifier}\Below{Name} \B[-5]{d} & & \K{Identifier}\Below{Token:}\Below{CategoryInfo} & & \K{Accessor}\Below{List} \B[-5]{dll} \B[-5]{d} \B[-5]{drr} \\
& & & & \K{Identifier}\Below{Token:}\Below{CategoryInfoResult} & & \K{OpenBrace}\Below{Token} & & \K{GetAccessor}\Below{Declaration} \B[-5]{dl} \B[-5]{dr} & & \K{CloseBrace}\Below{Token} \\
& & & & & & & \K{Get}\Below{Keyword} & & \K{Semicolon}\Below{Token}
}
\end{center}

\caption[Example of piece of C\# code and its Roslyn syntax tree]
{Example of piece of C\# code and its Roslyn syntax tree \\ (trivia not shown)}
\label{Syntax tree}
\end{figure}

Roslyn syntax trees are immutable and can be created using factory methods from the \lstinline{Syntax} class.
And while not all elements of the syntax tree have to be specified (like braces of a property accessor list),
creating a syntax tree can be quite cumbersome.

The exact syntax for creating syntax trees changed between the first two \acp{CTP}.
In the October 2011 \ac{CTP}, methods with many optional parameters were used.
In the June 2012 \ac{CTP}, the situation somewhat improved:
the factory method now has parameters only for required children of the created node
and optional child nodes can be added in a fluent manner using \lstinline{With*} methods.
The syntax used did not change in the September 2012 \ac{CTP}.

For an example of code to manually create the syntax tree from Figure~\ref{Syntax tree},
see Figure~\ref{Roslyn code 2011} for the October 2011 \ac{CTP} version and
Figure~\ref{Roslyn code 2012} for the September 2012 \ac{CTP} version.

\begin{figure}[p]

\begin{lstlisting}
Syntax.PropertyDeclaration(
  modifiers:
    Syntax.TokenList(
      Syntax.Token(SyntaxKind.PublicKeyword),
      Syntax.Token(SyntaxKind.AbstractKeyword)),
  type: Syntax.ParseTypeName("CategoryInfoResult"),
  identifier: Syntax.Identifier("CategoryInfo"),
  accessorList:
    Syntax.AccessorList(
      accessors:
        Syntax.List(
          Syntax.AccessorDeclaration(
            SyntaxKind.GetAccessorDeclaration,
            semicolonTokenOpt:
              Syntax.Token(SyntaxKind.SemicolonToken)))))
\end{lstlisting}

\caption{Sample code to manually create a Roslyn syntax tree \\ using October 2011 \ac{CTP}}
\label{Roslyn code 2011}
\end{figure}

\begin{figure}[p]

\begin{lstlisting}
Syntax.PropertyDeclaration(
  Syntax.ParseTypeName("CategoryInfoResult"),
  "CategoryInfo")
  .WithModifiers(
    Syntax.TokenList(
      Syntax.Token(SyntaxKind.PublicKeyword),
      Syntax.Token(SyntaxKind.AbstractKeyword)))
  .WithAccessorList(
    Syntax.AccessorList(
      Syntax.List(
        Syntax.AccessorDeclaration(
          SyntaxKind.GetAccessorDeclaration)
          .WithSemicolonToken(
            Syntax.Token(SyntaxKind.SemicolonToken)))))
\end{lstlisting}

\caption{Sample code to manually create a Roslyn syntax tree \\ using September 2012 \ac{CTP}}
\label{Roslyn code 2012}
\end{figure}

\chapter{MediaWiki improvements}
\label{mw improvements}

As mentioned in Section~\ref{paraminfo}, to generate types for each module of the \ac{API},
it is necessary to know the properties contained in the module response
and how do they map to the values of the \texttt{prop} parameter.

This information was not available, in MediaWiki previously.
For this reason, we extended the \texttt{paraminfo} module
to be able to provide information about result properties
of the \ac{API} modules, using the same type system already used to describe parameters.
Also, most of the \ac{API} modules were changed so that they provide this information to the \texttt{paraminfo} module.

Specifically, this meant adding \texttt{getResultProperties()} function to the \texttt{Api\-Base} class,
which is a base class for all module classes, and then overriding it in each module's class.
This function is then called from the \texttt{ApiParamInfo} class,
so that the provided information shows in the output of the \texttt{paraminfo} module.
For an example implementation of \texttt{getResultProperties()}, see Figure~\ref{paraminfo code}.

Of the 73~modules present in the MediaWiki core (that is, without any extensions),
5 are not suitable for having their result properties described,
because their result looks different than the result of other modules (for example, there are modules that produce \acs{RSS} feeds).
Further 5 modules do use the same response format as the other modules,
but their response cannot be described in the type system used.
There are also 17 modules that can be partially represented using this type system, but not completely.

The patch that adds this ability to the \texttt{paraminfo} module and the necessary
information to most other modules was reviewed by MediaWiki developers and merged into the official repository
on 12 June 2012.
On 2 July 2012, MediaWiki version 1.20wmf6, which includes changes from this patch, was deployed to all Wikimedia sites, including Wikipedias.

An example of the added result information to the \texttt{paraminfo} response (here for the \texttt{categorymembers} module) is in Figure~\ref{paraminfo props}.

\begin{figure}[htbp]
\begin{lstlisting}[language=php]
public function getResultProperties() {
  return array(
    'ids' => array(
      'pageid' => 'integer'
    ),
    'title' => array(
      'ns' => 'namespace',
      'title' => 'string'
    ),
    'sortkey' => array(
      'sortkey' => 'string'
    ),
    'sortkeyprefix' => array(
      'sortkeyprefix' => 'string'
    ),
    'type' => array(
      'type' => array(
        ApiBase::PROP_TYPE => array(
          'page',
          'subcat',
          'file'
        )
      )
    ),
    'timestamp' => array(
      'timestamp' => 'timestamp'
    )
  );
}
\end{lstlisting}

\caption{Implementation of \texttt{getResultProperties()} \\
in the \texttt{Api\-Query\-Category\-Members} class}
\label{paraminfo code}
\end{figure}

\begin{figure}[p]

\begin{lstlisting}[language=]
<props>
  <prop name="ids">
    <properties>
      <property name="pageid" type="integer" />
    </properties>
  </prop>
  <prop name="title">
    <properties>
      <property name="ns" type="namespace" />
      <property name="title" type="string" />
    </properties>
  </prop>
  <prop name="sortkey">
    <properties>
      <property name="sortkey" type="string" />
    </properties>
  </prop>
  <prop name="sortkeyprefix">
    <properties>
      <property name="sortkeyprefix" type="string" />
    </properties>
  </prop>
  <prop name="type">
    <properties>
      <property name="type">
      <type>
        <t>page</t>
        <t>subcat</t>
        <t>file</t>
      </type>
      </property>
    </properties>
  </prop>
  <prop name="timestamp">
    <properties>
      <property name="timestamp" type="timestamp" />
    </properties>
  </prop>
</props>
\end{lstlisting}

\caption{Result properties information for the \texttt{categorymembers} module}
\label{paraminfo props}
\end{figure}

During this work, we also noticed several bugs and inconsistencies in the \ac{API}.
Because of this, we reported eight bugs to the WikiMedia bug-tracking system.
Three of them turned out to be duplicates of already reported bugs and,
as of December 2012, three of them are still waiting to be fixed.

We also submitted eight additional patches to the MediaWiki code review system.
Although only three of them actually fix behavior of the MediaWiki \ac{API},
the rest are only fixes in documentation and other mostly insignificant changes.
Of those three patches, one is still waiting for review, because it is a breaking change,
and most likely will not be accepted.
\chapter{The LinqToWiki library}
\label{ltw}

The LinqToWiki library consists of one Visual Studio solution, that contains the following projects:

\begin{itemize}
\item LinqToWiki.Core
\item LinqToWiki.Codegen
\item LinqToWiki.Codegen.App
\item LinqToWiki.ManuallyGenerated
\item LinqToWiki.Samples
\end{itemize}

The LinqToWiki.Core project contains the core of the library:
types that access the \ac{API}, convert to and from the representation of data in the \ac{API},
represent parameters of various types of queries, represent query results
or those that process \ac{LINQ} expression trees.
This project can be used together with code generated using LinqToWiki.Codegen,
or with manually written code.

The LinqToWiki.Codegen project handles generating code based on information from the \texttt{paraminfo} module.
It contains types that represent the results of that module, process them, generate C\# code and compile this code.
This project also contains helper types for easier creating of Roslyn syntax trees.

The LinqToWiki.Codegen.App project compiles down to a simple console application called linqtowiki-codegen,
that uses functionality from the LinqToWiki.\allowbreak{}Codegen project.

The LinqToWiki.ManuallyGenerated project is a sample of how one could write code to access a wiki using LinqToWiki without using LinqToWiki.Codegen to generate the code.

Finally, the LinqToWiki.Samples project contains samples showing how to use various API modules using LinqToWiki.
It uses code generated by LinqToWiki.\allowbreak{}Codegen.\allowbreak{}App.

\paragraph{Usage}

The intended usage of LinqToWiki is this:
First run the linqtowiki-codegen application to generate a \ac{DLL} tailored for a certain wiki.
Then use the generated library together with LinqToWiki.Core in your C\# (or \ac{VB.NET}) application to access that wiki.

An alternative is to get the generated \ac{DLL} from someone else and then use that.
An advantage of this approach is that the user does not have to have Roslyn installed.

Other options are possible, though.
For example, the LinqToWiki.Codegen library can be used to generate the code as a set of files containing C\# source code.
Those files can then be modified and manually compiled.

\section{The LinqToWiki.Core project}

The LinqToWiki.Core project contains shared code that can be used when querying any MediaWiki wiki
that has the \ac{API} enabled.
It can be used together with code generated through LinqToWiki\allowbreak{}.Codegen,
but it can also be used without it.

In fact, LinqToWiki.\allowbreak{}Codegen internally uses LinqToWiki.Core to access the \texttt{paraminfo} module
using manually written code.

\subsection{\texorpdfstring{\lstinline{QueryTypeProperties}}{QueryTypeProperties}}

The \lstinline{QueryTypeProperties} class holds basic information about a “query type”,
which corresponds to an \ac{API} module.
This information includes the prefix this module uses in its parameters,
what type of module it is or mapping of its result properties to values accepted by the \texttt{prop} parameter.
It is also able to parse \ac{XML} elements this module returns.

\subsection{\texorpdfstring{\lstinline{WikiQuery}}{WikiQuery}}

Probably the most often used and certainly the most interesting queries are those using \texttt{list} query modules.
Such queries are represented in LinqToWiki by a group of types whose names start with \lstinline{WikiQuery}.

Specifically, there are four such types:
\lstinline{WikiQuery}, \lstinline{WikiQuerySortable}, \lstinline{Wiki}\lstBreak\lstinline{Query}\lstBreak\lstinline{Generator} and \lstinline{WikiQuerySortableGenerator}.
If a module supports sorting, it is represented by a type with \lstinline{Sortable} in its name
and if it supports being used as a generator for \texttt{prop} queries, it is represented by a type with \lstinline{Generator} in its name.

There is also a fifth type: \lstinline{WikiQueryResult}.
This type by itself represents a query that cannot be modified anymore,
but can be used to execute it and get the results.
All of the four preceding types inherit from \lstinline{WikiQueryResult},
so it is possible  to execute the query using any one of them too.

The type governs what operations are available.
For example, if a type is one of the two \lstinline{Sortable} types,
it will have an \lstinline{OrderBy()} method, but no other type has this method.
Each method can also return a different type, as is necessary to form queries.

\paragraph{\lstinline{WikiQuery} and generics}

All of the \lstinline{WikiQuery}-related types are generic
and their type parameters are used to decide what properties can be used in each operation.
For example, the type parameter \lstinline{TOrderBy} of \lstinline{WikiQuerySortable}
decides what properties can be used in the parameter of the \lstinline{OrderBy} method.

The way this is achieved is that \lstinline{TOrderBy} is a type that contains the properties that can be
used for sorting in the module \lstinline{WikiQuerySortable} represents
and the \lstinline{OrderBy} method accepts lambda expressions whose parameter is of this type.

For example, if some module supported sorting by \lstinline{PageId} and \lstinline{Title},
then \lstinline{TOrderBy} would be a type that contains two properties with those names.
Because of this, a query like \lstinline{source.OrderBy(x => x.Title)} would compile and execute fine,
but \lstinline{source.OrderBy(x => x.Name)} would fail to compile.

Because of the way lambda expressions work, queries like \lstinline[breaklines=true]{source.OrderBy(x => x.Title.Substring(1))} or \lstinline{source.OrderBy(x => random.Next())} would compile fine.
But because there is no way to efficiently execute such queries using the MediaWiki \ac{API},
they will fail with an exception at runtime.
This is a well-known problem with LINQ that also affects LINQ to SQL \cite{linq-to-sql-functions}
and other LINQ providers.

\paragraph{\lstinline{WikiQuery} operations}

The standard query operators available on the \lstinline{Wiki}\lstBreak\lstinline{Query} types are:

\begin{itemize}
\item \lstinline{Where()} only sets some parameter or parameters of a query,
it always returns the same type.

It is available on all four of the basic \lstinline{WikiQuery} types
and uses the generic type parameter \lstinline{TWhere}.

\item \lstinline{Select()} is used to choose how the elements in the resulting collection should look like
and what properties should they contain.
Because the result of the lambda passed into this method can be an arbitrary type,
it doesn't make sense to modify the query after calling this method.
Because of that, \lstinline{Select()} returns \lstinline{WikiQueryResult}.
This also follows query expression syntax, where \lstinline{select} is the last clause of each query.

It is available on all four of the \lstinline{WikiQuery} types
and uses the type parameter \lstinline{TSelect}.

\item \lstinline{ToEnumerable()} and \lstinline{ToList()} are used to actually execute the query.
The distinction between the two methods is that \lstinline{ToEnumerable()} returns an \lstinline{IEnumerable},
that lazily loads new pages of results on demand.
\lstinline{ToList()}, on the other hand, returns a \lstinline{List},
that is immediately loaded with all of the results, possibly from many pages.

These two methods are available on all of the \lstinline{WikiQuery} types, including \lstinline{WikiQueryResult}
and return the result based on the type parameter \lstinline{TSource} for most of the types.
An exception is \lstinline{WikiQueryResult}, which uses a separate \lstinline{TResult} type parameter.

\item \lstinline{OrderBy()} (and \lstinline{OrderByDescending()}) sets the ordering.
Because it does not make sense to sort the same query multiple times
and because no module supports sorting by multiple keys,
this method returns the type with \lstinline{Sortable} removed.

This method is available on the two \lstinline{Sortable} types
and uses the type parameter \lstinline{TOrderBy}.

\item \lstinline{Pages} is a property that returns a \lstinline{PagesSource}
that can then be used in a \texttt{prop} query.
See Section~\ref{PagesSource} for more information.

This property is available on the two \lstinline{Generator} types
and uses the type parameter \lstinline{TPage}.

\end{itemize}

For a state diagram of transitions between the \lstinline{WikiQuery} types and other related types,
see Figure~\ref{WikiQuery types}.

\begin{figure}[htbp]
\begin{center}
\begin{tikzpicture}[>=angle 90]
\path (0,9) node(WQS) {\lstinline{WQSortable}};
\path (5,9) node(WQ) {\lstinline{WQ}};

\path (5,6) node(WQR) {\lstinline{WQResult}};

\path (0,3) node(WQSG) {\lstinline{WQSortableGenerator}};
\path (5,3) node(WQG) {\lstinline{WQGenerator}};

\path (0,0) node(PS) {\lstinline{PagesSource}};
\path (5,0) node(WQPR) {\lstinline{WQPageResult}};

\draw[->] (WQS) edge [out=188,in=172,looseness=5,auto] node {\lstinline{Where}} (WQS);
\draw[->] (WQS) edge node[above] {\lstinline{OrderBy}} (WQ);
\draw[->] (WQS) edge node[below left] {\lstinline{Select}} (WQR);

\draw[->] (WQ) edge [out=-20,in=20,looseness=8,right] node {\lstinline{Where}} (WQ);
\draw[->] (WQ) edge node[right] {\lstinline{Select}} (WQR);

\draw[->] (WQSG) edge [out=185,in=175,looseness=5,auto] node {\lstinline{Where}} (WQSG);
\draw[->] (WQSG) edge node[auto] {\lstinline{OrderBy}} (WQG);
\draw[->] (WQSG) edge node[above left] {\lstinline{Select}} (WQR);

\draw[->] (WQG) edge [out=-8,in=8,looseness=5,right] node {\lstinline{Where}} (WQG);
\draw[->] (WQG) edge node[right] {\lstinline{Select}} (WQR);

\draw[->] (WQSG) edge node[auto] {\lstinline{Pages}} (PS);
\draw[->] (WQG) edge node[auto] {\lstinline{Pages}} (PS);

\draw[->] (PS) edge node[auto] {\lstinline{Select}} (WQPR);

\end{tikzpicture}
\end{center}

\caption[State diagram of \lstinline{WikiQuery}-related types]
{State diagram of \lstinline{WikiQuery}-related types \\ (\lstinline{WikiQuery} is shortened to \lstinline{WQ} to save space)}
\label{WikiQuery types}

\end{figure}

\subsection{\texorpdfstring{\lstinline{PagesSource}}{PagesSource}}
\label{PagesSource}

The \lstinline{PagesSource} type represents a collection of pages that can be used in \texttt{prop} queries,
to get information about those pages.
This information can be for example a list of categories for each page in the collection.

There are two kinds of \lstinline{PagesSource}s: generator-based and list-based:

\begin{itemize}

\item List-based sources use a static list of pages, given as a collection of page titles, page IDs or revision IDs.

Because the number of pages given this way in a single \ac{API} request is fairly limited (usually to 50),
large lists have to be queried multiple times.
\lstinline{PagesSource} handles this transparently, so the user can input as many pages as he wants and does not have to worry about the limit.

One exception is if the limit is different than the default of 50 for the current user on the current wiki.
In that case, the user should change the limit by setting the property \lstinline{PagesSourcePageSize} on the \lstinline{Wiki} object.%
\footnote{In all other cases where limits are important in this library, they limit the output, not the input.
That is why simply setting \texttt{limit=max} works in those other cases, but does not work here.}

If the collection used to create a \lstinline{PagesSource} is lazy, it is iterated in a lazy manner.
For example, it could be the result of another LinqToWiki query, with additional processing by \ac{LINQ} to Objects,
that is not possible using LinqToWiki alone.
Or it could the result of a query from another wiki.
In such cases, the original query will only make as many requests as necessary for the follow-up query.

\item Generator-based sources represent a dynamic list of pages that is the result of another \ac{API} query,
like the list of all pages on a wiki from the \texttt{allpages} module.
This way, the list of pages does not have to be retrieved separately, only to be sent back.

Generator queries also have to handle paging, as described in Section~\ref{mw paging},
including the exception for the \texttt{revisions} module.

\end{itemize}

Thanks to the fact that both kinds of page sources for \texttt{prop} queries are represented by the same
(abstract) type, the user of this library can use the same code to work with any source,
thus avoiding repetitive code.

\paragraph{Structure of \texttt{prop} query}

To actually create a \texttt{prop} query for a page source, one uses the \lstinline{Select()} method.
Its parameter is a lambda, whose parameter is the type parameter \lstinline{TPage} of \lstinline{PagesSource}.
This type is the same for all queries on the same wiki, but could be different for different wikis.

Inside the lambda, properties and methods of the \lstinline{TPage} type can be accessed.
Each of them represents a \texttt{prop} module and all of the methods return one of the \lstinline{WikiQuery} types,
which can then be queried as usual, with one condition:
the \lstinline{WikiQuery} types can't ``leak'' outside of the query, so one has to use \lstinline{ToEnumerable()} or \lstinline{ToList()} inside the lambda.

There is a special case for the \texttt{revisions} module:
it can be also used with the \lstinline{FirstOrDefault()} method,
which means only the most recent revision for each page is selected.

If a \texttt{prop} module has a single result (not a collection), it is represented as a property
that directly returns this result, no querying is possible.

The methods of these \texttt{prop} queries are inside a lambda expression,
so they are not actually executed unless the expression was compiled and the resulting delegate invoked.
Because of this, processing them is not as simple as with normal queries.
For more details, see Section~\ref{PageExpressionParser}.

\medskip

For an example of \lstinline{PagesSource} query, see Figure~\ref{PS query}.

\begin{figure}[htbp]

\begin{lstlisting}
pagesSource.Select(
    p =>
    new
    {
        p.Info,
        Categories =
            p.Categories()
            .Where(c => c.Show == Show.NotHidden)
            .Select(c => new { c.Title, c.SortKeyPrefix })
            .ToEnumerable()
            .Take(10)
	}
)
\end{lstlisting}

\caption{Example of \lstinline{PagesSource} query \\ that uses \texttt{info} and \texttt{categories} modules}
\label{PS query}

\end{figure}

\needspace{5\baselineskip}

\subsection{\texorpdfstring{\lstinline{QueryParameters}}{QueryParameters}}

The \lstinline{QueryParameters} type contains the parameters of a query:

\begin{itemize}
\item sort direction and parameter by which to sort,
\item list of properties to select and a delegate that uses them to construct the result object,
\item list of other parameters, as key-value pairs.
\end{itemize}

\lstinline{QueryParameters} is an immutable type,
so that an initial subquery can be safely used repeatedly, as is the case with LINQ to Objects.
The list of other parameters is a functional-style immutable linked list.

\medskip

The \lstinline{PropQueryParameters} type derives from \lstinline{QueryParameters}
and is used to store information about a single module in a \texttt{prop} query.
Apart from inherited members, it also contains the name of the module
and a special value indicating whether to retrieve only the first item,
which corresponds to the usage of the \lstinline{FirstOrDefault()} method.

A related type is \lstinline{PageQueryParameters},
which represents a whole \texttt{prop} query.
That means it contains a list of \lstinline{PropQueryParameters} objects
and also information about the source of the query.

\subsection{\texorpdfstring{\lstinline{ExpressionParser}}{ExpressionParser}}

The \lstinline{ExpressionParser} static class is used to process expression trees from \ac{LINQ} methods
and store the processed query parameters in \lstinline{QueryParameters}.

Common for all expression tree processing is that closed-over local variables contained in the processed lambda,
which are represented as members of a compiler-generated closure class,
have to be first replaced by their actual value.
This is done using \lstinline{PartialEvaluator} written by Matt Warren \cite{warren}.

Also, some property names have to be translated from their C\# version to their \ac{API} version.
For details and the reason why this is necessary, see Section~\ref{ltwc-naming}.

\medskip

Each of the methods requires different processing. Specifically:

\begin{itemize}
\item Expression trees from \lstinline{Where()} are first split into one or more subexpressions
that are \emph{and}ed together (\lstinline{x => subexpr1 && subexpr2 && ...}; \emph{or} is not supported by the \ac{API})
and each of the subexpressions is then added as a key-value pair to the result.

Each subexpression has to be in the form \lstinline{x.Property == Value},
where \lstinline{Value} is a constant, possibly from an evaluated closed-over variable.
The reverse order (\lstinline{Value == obj.Property}) is also allowed.
An alternative for boolean properties is accessing the property directly (\lstinline{x.Property})
or negated (\lstinline{!x.Property}).

\item Processing \lstinline{OrderBy()} expression trees is simple:
they can either be identities (\lstinline{x => x}), which means default sorting will be used
(which is the only possibility for some modules),
or they can be simple property accesses (\lstinline{x => x.Property}),
which means the result will be sorted by that property.

The order of sorting (ascending or descending) is decided by the method used:
whether it was \lstinline{OrderBy()} or \lstinline{OrderByDescending()}.

\item Expression trees from \lstinline{Select()} are processed in two steps.
First, the expression is scanned for usages of its parameter.
If any of its properties are used, it means those properties have to be retrieved from the \ac{API}.
If the parameter is used directly, without accessing its properties,
it means all of the properties have to be retrieved, because it is impossible to say which of them will be used.

For example, the expression \lstinline|x => new { x.Property1, x.Property2 }|
means only \lstinline{Property1} and \lstinline{Property2} have to be retrieved.
On the other hand, \lstinline{x => SomeMethod(x)} means all of the properties have to be retrieved.

Second step is compiling the expression into a delegate,
which will then be executed for each item coming from the \ac{API}.

Put together, these two steps mean that \lstinline{Select()} can be used with any expression
and only properties that are actually needed will be returned by the \ac{API}.
\end{itemize}

\subsection{\texorpdfstring{\lstinline{PageExpressionParser}}{PageExpressionParser}}
\label{PageExpressionParser}

The class \lstinline{PageExpressionParser} is used to process
the \lstinline{Select()} lambda in \lstinline{PagesSource} queries.
The difficulty there is that the direct approach of building the query step-by-step,
used in normal queries, will not work.
That is because the expression has to be analyzed before there is any page object
that it expects as its parameter.

The result of this analysis is twofold:
the set of parameters needed for all of the \lstinline{prop} queries,
as a collection of \lstinline{PropQueryParameters},
and a delegate that can be used to get the result object for each page in the \ac{API} response.

\medskip

Because the subquery for each \texttt{prop} module has to end with a call to \lstinline{To}\lstBreak\lstinline{Enumerable()} or \lstinline{ToList()},
the parameters can be extracted by invoking the part of the subquery before that call.
At the beginning of each subquery is invoking a module-specific method on the page object.
But because there is no page object to use, that invocation is first replaced by an appropriate \lstinline{WikiQuery} object.

For example, for the query in Figure~\ref{PS query}, the invoked code is (where \lstinline{wikiQuery} is the appropriate \lstinline{WikiQuery} object):

\begin{lstlisting}
wikiQuery.Where(c => c.Show == Show.NotHidden)
	  .Select(c => new { c.Title, c.SortKeyPrefix })
\end{lstlisting}

\medskip

To get the delegate, all calls to \lstinline{Where()} and \lstinline{OrderBy()} are removed,
because their only purpose is to modify the query parameters.
Then the single parameter of type \lstinline{TPage} is replaced by a parameter of type \lstinline{PageData}
and calls to module methods are replaced by calls to \lstinline{GetData()},
with a type parameter specifying the type of the result and a parameter specifying the name of the module.

The \lstinline{GetData()} method returns a collection,
so for modules that return only a single item, like \texttt{info},
a call to \lstinline{Single}\lstBreak\lstinline{Or}\lstBreak\lstinline{Default()} is also added.

\needspace{11\baselineskip}

For example the expression in the query in Figure~\ref{PS query} is transformed into:

\begin{lstlisting}
pageData =>
new
{
    Info = pageData.GetData<InfoResult>("info")
    		     .SingleOrDefault(),
    Categories =
    	pageData.GetData<CategoriesSelect>("categories")
    		 .Select(c => new { c.Title, c.SortKeyPrefix })
    		 .Take(10)
}
\end{lstlisting}

\subsection{Other types}

The \lstinline{QueryProcessor} type manages downloading the result and transforming it from \ac{XML} to objects.
For queries whose result is a collection, it also handles returning the pages in a lazy manner
and downloading the follow-up pages when necessary.

The \lstinline{QueryPageProcessor} type does the same for \lstinline{PagesSource} queries.

\medskip

The \lstinline{Downloader} type takes care of forming the query string, executing the request and
returning the result as an \lstinline{XDocument}.
\lstinline{XDocument} is a part of LINQ to XML, a part of .NET Framework for manipulating \ac{XML} documents.

\lstinline{Downloader} always uses POST and formats its requests as \path{application/x-www-form-urlencoded}.
This means that all modules work, including those that require POST.
On the other hand, uploads of files do not work, because they require \path{multipart/form-data}.

The decision to use \path{application/x-www-form-urlencoded} follows from the fact that
\path{multipart/form-data} is very inefficient when sending multiple parameters with short values,
which is common when making requests to the \ac{API}.

% TODO: Namespace type?

\section{The LinqToWiki.Codegen project}
\label{ltwc}

The LinqToWiki.Codegen project contains code that retrieves information about \ac{API} modules in some wiki,
then uses that information to generate C\# code to access those modules using Roslyn
and finally compiles the code into a library.

\paragraph{Alternatives to Roslyn}

Roslyn was chosen, because it is superior when compared with common approaches for code generation in .NET,
namely Reflection.Emit and \acs{CodeDOM}.

Reflection.Emit \cite{reflection-emit} is a set of types that allow generation of code at runtime.
The generated code can then be directly executed or saved as an assembly (.dll or .exe) to disk.
The distinguishing feature is that it uses the low-level \ac{CIL},
which means writing any code beyond the simplest methods can be very tedious and error-prone.

\acs{CodeDOM} \cite{codedom} can be used to generate code and compile it to an assembly.
It uses language-independent model, which can be converted to various .NET languages,
including C\# and \ac{VB.NET}.
This model is also the biggest disadvantage of \acs{CodeDOM}, because it means it doesn't support all features of C\#.
For example, even such basic feature as writing a \lstinline{static} class is impossible in the \acs{CodeDOM} model
without using “hacks”.

Detailed description of Roslyn is in Section~\ref{roslyn}.

\paragraph{What to generate?}

At this point, we have a library (LinqToWiki.Core) that can be used to access the MediaWiki \ac{API} the way we want
from the final generated library.
We can also use the same library to get the information we need about the modules of the \ac{API} from the \texttt{paraminfo} module.
And we have decided we want to use Roslyn to generate the final library.
What remains is to decide what code to generate, how exactly to map the modules, their parameters
and their results into the model of LinqToWiki.Core.

%TODO: think about moving this before LTW.Core

There are some decisions that were already made in LinqToWiki.Core
(the \texttt{sort} and \texttt{dir} parameters should map to \lstinline{OrderBy()};
the \texttt{prop} parameter maps to \lstinline{Select()}),
but several other decisions still remain:\footnote{
Obviously, both libraries were written alongside each other, to work well together, not one after the other.
But we think it is better to describe them this way, separately.}

\begin{itemize}
\item How should the remaining parameters be mapped?
Should they all go into \lstinline{Where()} or somewhere else? Where?
\item How should the modules that do not return lists be mapped?
\ac{LINQ} methods are not suitable for them, because they are meant to work with collections.
\item How to name the generated types and members?
Specifically, how to represent names that cannot be used (like those containing special characters)
and names that are undesirable (those that conflict with C\# keywords).
Also, should the generated members follow .NET naming conventions?
\end{itemize}

Our answers to these questions are in the following couple of sections.

\subsection{Naming of generated types and members}
\label{ltwc-naming}

Let us start with the last question: Should the generated members follow .NET naming conventions?
The .NET naming guidelines \cite{guidelines-for-names},
that are widely followed by various .NET libraries and the .NET Framework itself,
state that names of types and public members should use PascalCase,
that is, each word of an identifier should start with a capital letter
and the identifier should not contain any delimiters (such as underscores).

We would prefer to follow these naming conventions, but, unfortunately, it is not possible.
That is because the names of modules, parameters, result properties and almost all enumerated types in the \ac{API}
use names that are all lowercase, without delimiters between words.
That means there is no way to figure out which letters in an identifier should be capitalized
(apart from the first one).

As one of the more extreme examples,
one of the possible values of the \texttt{rights} parameter
of the \texttt{allusers} module on the English Wikipedia is \texttt{collection\-save\-as\-community\-page}.
A human can see that the proper name for that value using PascalCase would be \lstinline{CollectionSaveAsCommunityPage},
but a computer cannot.
(Actually, it is possible that the words could be reliably separated using natural language processing,
but doing that is outside the scope of this work.)

\paragraph{Reserved names}

Because different .NET languages have different sets of reserved identifier names
(usually, those are the language keywords)
and because libraries written in one language should be usable from other languages,
.NET languages provide a way to use their keywords as identifiers.
In the case of C\#, this is done by prefixing the identifier with an at sign.
So, for example, to use \texttt{new} as an identifier, one has to write \lstinline{@new}.

Thanks to this, using keyword-named identifiers is still possible,
although slightly less convenient than with normal identifiers.
Also, the naming guidelines suggest avoiding keywords as identifiers.

In MediaWiki core \ac{API} modules, there are four identifiers that are also C\# keywords:
\texttt{namespace}, \texttt{new}, \texttt{true}, \texttt{false}.
Out of these, we decided to shorten \texttt{namespace} to \lstinline{ns},
which is a common abbreviation, so the meaning should not be lost.
The other three have to be written with \lstinline{@} (\lstinline{@new}, \lstinline{@true} and \lstinline{@false}) in C\#,
because we did not find a reasonable alternative for them.

\paragraph{Special characters}

The characters hyphen~(\texttt{-}), slash~(\texttt{/}) and space
appear in some names in the \ac{API}, but are not allowed in .NET identifiers,
so they are replaced by underscores (\texttt{\_}).

Some names also start with an exclamation mark~(\texttt{!}), to indicate negation.
Such names are translated by prefixing \lstinline{not_}.
So for example, \texttt{!minor} (which means that an edit is not a minor edit)
is translated into \lstinline{not_minor}.

One more special case is that some enumerated types allow an empty value.
Such value is then represented by the identifier \lstinline{none}.

\paragraph{Naming types}

Another question is how to name the generated types.
There are two kinds of generated types:
those that represent some enumerated type and those that represent parameters or results of some module.

For the latter kind, it is simple to come up with a convention like naming them by the module name,
suffixed by the specific kind of the type
(e.g. \lstinline{blockResult} for the result of the \texttt{block} module
or \lstinline{categorymembersWhere} for the type representing \lstinline{Where()} parameters for the \texttt{categorymembers} module).

But for the former kind, the situation is more complicated.
Enumerated types do not have names by themselves, they are part of a parameter or property that has a name.
The problem is that different modules often have parameters and properties with the same name,
while their type sometimes is the same and sometimes it is not.

So, there are two options: either let the types that look the same actually be the same generated type,
or let each parameter and property have its own distinct type.
If we merge the types that look the same, we should not use the module name in their name,
because one type can be used with different modules.
But that means we need to distinguish different types in another way, like a number.
But names like \lstinline{token5} are not very helpful for the user.

Because of that, we chose the other option, which means including the name of the module in the type name.
But doing it this way does not eliminate conflicts completely:
In the case when a module has a parameter and a property with the same name,
their types still have to be distinguished.
An example of such type name is \lstinline{recentchangestype2}.

\subsection{Structure of generated code}
\label{ltwcg-structure}

At the start of each query is the \lstinline{Wiki} type.
It contains methods for non-query modules as well as methods to create list-based \lstinline{PagesSource}s.
It also contains the property \lstinline{Query} that returns an object that contains methods for
\texttt{list} and \texttt{meta} query modules.
(\texttt{prop} query modules work differently, for more information, see Section~\ref{PagesSource}.)

\paragraph{Simple modules}

With modules that do not return lists, the situation is mostly simple:
there are no parameters to sort or filter the result (because it is not a list)
and most of those modules also do not have parameters to choose the result properties.

Because of that, a method for each such module, that directly returns the result object is enough.
This method has parameters corresponding to the parameters of the module,
where required parameters of the module are mapped as normal method parameters
and parameters that are not required are mapped as optional parameters.
The code of this method builds \lstinline{QueryParameters} from the method parameters
and then executes the query using \lstinline{QueryProcessor}.

\paragraph{List modules}

On the other hand, list modules can have several kinds of parameters:

\begin{itemize}
\item Those that affect order of the items in the list. They are naturally mapped as \lstinline{OrderBy()}.
The parameters \texttt{sort} and \texttt{dir} belong here.
\item Those that choose what properties appear in the result. They are naturally mapped as \lstinline{Select()}.
Only the parameter \texttt{prop} belongs here.
\item Those that filter what items appear in the result. They are naturally mapped as \lstinline{Where()}.
For example, the parameters \texttt{namespace} and \texttt{startsortkey}
of the \texttt{categorymembers} module belong here.
\item Various other parameters. They do not naturally map to any LINQ method.
For example, the parameter \texttt{title} (that decides which category to enumerate)
of the \texttt{categorymembers} module belongs here.
\end{itemize}

The first two kinds are not a problem, because it is clear which parameters belong to them.
The second two kinds are a problem, because there is no clear way to automatically distinguish between the two.
One exception, where the distinction is clear,
is if a parameter is required (as indicated in its description),
because it then means it belongs among “other parameters”.

Required parameters are given as parameters of the module methods,
but we decided to treat all non-required parameters that do not belong to the first two kinds,
as if they were \lstinline{Where()} parameters.
Unfortunately, this means that some queries do not logically make sense,
if we consider that the \lstinline{Where()} method should only filter the results.

For example, consider this query:

\begin{lstlisting}
wiki.Query.categorymembers()
    .Where(cm => cm.title == "Category:Query languages")
\end{lstlisting}

Here, the \lstinline{title} property does not actually represent filtering
by the title of the category member, it decides which category to enumerate.
And without it, the query would not even execute successfully
(the parameter \texttt{title} is not marked as required, because the parameter \texttt{pageid} can be used instead of it).

Proper solution to this problem would require human interaction when generating the code,
to choose which parameters belong to \lstinline{Where()} and which do not.
As an alternative, the description of each parameter in the \texttt{paraminfo} module could contain its kind.

\paragraph{Enumerated types}

One more question is how to represent enumerated types.
The answer is seemingly simple: make them \lstinline{enum}s
and for those parameters or properties that can have multiple values, use bit flags.
But the largest type that can be used as an underlying type for \lstinline{enum}
is \lstinline{ulong}, which has 64 bits.
That means this will work only if there is no enumerated type in the \ac{API},
that has more than 64 values and can have multiple values at the same time.
Unfortunately, the English Wikipedia has one:
the type of the \lstinline{rights} parameter of the \lstinline{allusers} module
has 106 values and the parameter can have multiple values at the same time.

Because of that, each enumerated type is represented by immutable class deriving from the common base class \lstinline{StringValue},
with inaccessible constructor and static field for each possible value.
Combination of values can be represented as a collection, like with other types.

\subsection{\texorpdfstring{\lstinline{Wiki}}{Wiki}}

The top-level type that manages all code generation is \lstinline{Wiki}
(not to be confused with the generated \lstinline{Wiki} type from Section~\ref{ltwcg-structure}).
It manages retrieving information about \ac{API} modules and generating code for them.

When the code generation is complete, it saves the generated C\# files to a temporary directory
and compiles them using CodeDOM.
CodeDOM is used for the compilation,
because its compiler is the full C\# compiler and can handle all features of C\# (unlike the CodeDOM object model).
The Roslyn compiler is not able to compile some useful expressions, such as collection initializers
(but the object model of Roslyn is complete).

\subsection{\texorpdfstring{\lstinline{ModuleSource}}{ModuleSource}}

The \lstinline{ModuleSource} class is used to retrieve information about modules of the \ac{API}
and transform it from \acs{XML} to objects, like \lstinline{Module}, \lstinline{Parameter} and \lstinline{Parameter}\lstBreak\lstinline{Type}.
This information comes from the \texttt{paraminfo} module and is fetched using LinqToWiki.Core.

In fact, this code can be viewed as a sample on how to use LinqToWiki.Core
without code generated by LinqToWiki.Codegen.
Generated code cannot be used to work with the \texttt{paraminfo} module,
because it is one of the modules, whose response is complicated
and does not fit into the simple type system used by \texttt{paraminfo}.

Because the addition of result properties to \texttt{paraminfo} was made as a part of this work
(see Section~\ref{mw improvements}) and so is quite recent, there is also another option to get this information:
\lstinline{ModuleSource} can accept a “props defaults” file, that contains the necessary information.
The file looks the same as \texttt{paraminfo} response (in XML format),
except it contains only the added information.
This file can be created from another wiki that can already provide this information,
or it can be written by hand.
It can be also useful to work with modules from extensions, that currently don't provide this information.

\subsection{\texorpdfstring{\lstinline{ModuleGenerator}}{ModuleGenerator}}

\lstinline{ModuleGenerator} and related types are the ones that actually generate code for each module using Roslyn.
Each type generates code for a certain kind of module,
so for example \lstinline{ModuleGenerator} works with non-query modules,
while \lstinline{QueryModuleGenerator} works with most query modules.

Each generator creates all the code that is necessary for that module.
For example, for a \texttt{list} query module,
this includes generating \lstinline{Where}, \lstinline{Select} and possibly \lstinline{OrderBy} classes,
method in the \lstinline{QueryAction} class
(which is returned by the \lstinline{Query} property of the \lstinline{Wiki} class)
and types for all its enumerated types.

Each of the generated types and methods also has \acs{XML} documentation comment attached,
based on description from \texttt{paraminfo}.
This means that a user of this library does not have to guess what each method or property means,
his \ac{IDE} can show him description for it.

These descriptions sometimes contain references to details of the \ac{API} that this library abstracts away.
For example, the description for the \texttt{unique} parameter of the \lstinline{alllinks} module says:
“Only show unique links. Cannot be used with generator or alprop=ids.”
The reference to \texttt{alprop} makes sense to someone who uses the \ac{API} directly,
but would be very confusing for a user of LinqToWiki.
Not only does LinqToWiki abstract away module prefixes (\texttt{al}),
it also doesn't expose the \texttt{prop} parameter directly
(the \lstinline{Select()} method is used instead).

\subsection{\texorpdfstring{\lstinline{SyntaxEx}}{SyntaxEx}}

As described in Section~\ref{roslyn}, creating Roslyn syntax trees can be cumbersome.
The \lstinline{SyntaxEx} class makes doing that easier by adding simpler alternatives
to the factory methods in Roslyn's \lstinline{Syntax} class.
The \lstinline{SyntaxEx} methods do not handle more complex cases,
so for those, using \lstinline{Syntax} is still necessary.

For example of how the code from Figure~\ref{Roslyn code 2012} can be written using \lstinline{SyntaxEx},
see Figure~\ref{SyntaxEx code}.

\begin{figure}[htbp]

\begin{lstlisting}
SyntaxEx.AutoPropertyDeclaration(
  new[]
  {
    SyntaxKind.PublicKeyword,
    SyntaxKind.AbstractKeyword
  },
  "CategoryInfoResult",
  "CategoryInfo",
  setModifier: SyntaxKind.PrivateKeyword,
  isAbstract: true)
\end{lstlisting}

\caption{Sample code to manually create \\ a Roslyn syntax tree using \lstinline{SyntaxEx}}
\label{SyntaxEx code}
\end{figure}

\medskip

Another improvement is that syntax nodes that represent declaration of property, field, parameter or variable
can be used to refer to them in later code, for example when assigning the value of a parameter to a property.
This is achieved by using implicit conversions and a helper type \lstinline{NamedNode}.
In Roslyn without this extension, it is necessary to extract the name of the syntax node
and use that to create \lstinline{IdentifierNameSyntax}.

As with \lstinline{SyntaxEx}, this can make simple cases simpler, but cannot handle everything.
Because of that, complex cases still have to directly use Roslyn.

\section{The \texorpdfstring{\lstinline{linqtowiki-codegen}}{linqtowiki-codegen} application}
\label{ltw-ca}

\lstinline{linqtowiki-codegen} is a simple console application
that can be used to access the functionality of LinqToWiki.Codegen.
In other words, it can generate a library for accessing a specific wiki using LinqToWiki.

Using its command-line arguments, one can specify the \ac{URL} of the wiki,
the directory to where the files will be generated, the namespace of the generated types,
the name of the generated assembly and the location of the props default file.

Some of the more advanced features of LinqToWiki.Codegen,
such as writing the generated C\# code to a specific directory,
are not available from this application.

The application writes a short usage note when run without arguments, which can be seen in Figure~\ref{lca-usage}.

\begin{figure}[htbp]

\begin{verbatim}
Usage:    linqtowiki-codegen url-to-api [namespace [output-name]]
[-d output-directory] [-p props-file-path]
Examples: linqtowiki-codegen en.wikipedia.org LinqToWiki.Enwiki
linqtowiki-enwiki -d C:\Temp -p props-defaults-sample.xml
          linqtowiki-codegen https://en.wikipedia.org/w/api.php
\end{verbatim}

\caption{Usage note of the \lstinline{linqtowiki-codegen} application}
\label{lca-usage}

\end{figure}

\section{Samples of queries}
\label{ltw-s}

The project LinqToWiki.Samples contains one class with methods
that show the usage of each available module of the \ac{API}.
It also contains one real-world complex query that combines LinqToWiki with LINQ to Objects
to search for empty categories that are not redirects.

When run, it is a console application that shows the output of the selected module.
Selecting the module to use is done by changing which method is called by the \lstinline{Main()} method.
\chapter{Future work}
\label{future work}

%TODO: mention < and >?

While the LinqToWiki library is fully functioning and could be considered complete,
there is still room for improvement.
Some possible improvements were already mentioned
(proper capitalization in generated code;
distinguishing \lstinline{Where()} parameters from other parameters),
but there are also other improvements that could have larger impact on the usage of the library:

\begin{itemize}
\item \emph{F\# implementation}

The recent version 3.0 of the .NET-based functional language F\# has its own alternative to \ac{LINQ},
called query expressions.
Query expressions use their own set of types, so they are incompatible with C\# \ac{LINQ},
but they offer similar capabilities.
Unlike C\#, they also support creating custom operators,
which could provide better syntax for \lstinline{PagesSource} queries.

Also, F\# 3.0 supports type providers, which is a feature for automatic creating of types
just before they are used in code, that is, before compilation.
This would simplify the workflow of using LinqToWiki, because the code generation would be automatic.

\item \emph{Asynchrony}

The current version of LinqToWiki is completely synchronous,
which means there is always a thread blocked, while an application waits for a network response.
Better support for asynchrony is the main feature of the new C\#~5.0,
but LinqToWiki could be made asynchronous even without it.
This could be useful especially in \ac{GUI} applications, where the main thread should not be blocked for long periods of time.

A big part of LinqToWiki is working with modules that return lists,
but there is no single idiomatic way of representing asynchronous collections,
with or without C\# 5.0 improvements.
There are several possibilities, including a lazy collection of \lstinline{Task}s,
\acs{Rx} \lstinline{IObservable} or \acs{TPL} Dataflow \lstinline{ISourceBlock}.

Another way how asynchrony could be used in LinqToWiki is if the following page was being retrieved
even before the user finished processing the preceding page.
This could be especially useful for \lstinline{PagesSource} queries,
because the following primary and secondary page could be fetched in parallel.

\item \emph{Better compile-time checking}

One of the features of Roslyn that is not currently used by LinqToWiki is Visual Studio integration.
What that means is that one can use Roslyn to write various Visual Studio extensions,
including enhanced compile-time checking.

In the case of LinqToWiki, this could be used to verify that queries that compile are actually correct.
If not, a custom error would show in the list of errors.
Such checking would be also useful for other \acs{LINQ} providers, such as LINQ to SQL,
but each provider has its own rules about what is an error and what is a correct query.

% extensions?

\end{itemize}
\chapter{Related work}
\label{related}

Many libraries for accessing MediaWiki written in various languages already exist.
Some relevant examples are included below.

\begin{itemize}
\item wikitools \cite{wikitools}

wikitools is a library written in Python. It uses string-based dictionaries for parameters and results.
For example, the \texttt{blocks} sample query used in Chapter~\ref{goal} would look like this:

\begin{lstlisting}[language=python]
params = {
    'action':'query',
    'list':'blocks',
    'bkip':'8.8.8.8',
    'bkdir':'older',
    'bkprop':'byid'
}
request = api.APIRequest(site, params)
request.query()
\end{lstlisting}

This is basically identical to the string-based approach mentioned in Section~\ref{alternatives},
so it also shares all of its disadvantages.

\item WikiFunctions \cite{wikifunctions}

WikiFunctions is a .NET library used primarily in a semi-automated MediaWiki editor AutoWikiBrowser.
It uses the manual approach for a small subset of available modules.
This means the user does not have to know much about the API,
but on the other hand, the user cannot use all of the functionality of the API.

\item Linq to Wikipedia \cite{linq-to-wikipedia}

Linq to Wikipedia is a .NET library that contains a simple \lstinline{IQueryable} provider for two modules
(\texttt{search} and \texttt{opensearch}) from the MediaWiki API.
A query using this library looks like this:

\begin{lstlisting}
from wikipedia in datacontext.OpenSearch
where wikipedia.Keyword == "linq"
select wikipedia
\end{lstlisting}

This query provider is very limited in that it supports
only \lstinline{Where()}, \lstinline{Take()} and \lstinline{Skip()} methods
(the \lstinline{select} clause in the above code is actually not compiled into a call to the \lstinline{Select()} method).
Other LINQ methods are attempted to be ignored and often cause exceptions at runtime.

Also, the \lstinline{where} clause has to be in a very specific shape:
it can use only the \lstinline{Keyword} property, even though the query object has other properties.
On the other hand, objects in the resulting collection have the \lstinline{Keyword} property,
even though it is always \lstinline{null}.
This is because with \lstinline{IQueryable}, the type of the result is the same as the type used in \lstinline{where}.

\end{itemize}
\chapwithtoc{Conclusion}
\label{conclusion}

The goal of this work was to implement a C\# library to access the MediaWiki \ac{API}
in a way that is readable, discoverable, strongly-typed and flexible.

These goals were successfully accomplished using a custom \acs{LINQ} provider and code generation with Roslyn \ac{CTP}:

\begin{itemize}
\item Code written using LinqToWiki is readable:
non-query modules are accessed using simple methods,
with parameters usually given as named parameters of the method.

Query modules use \acs{LINQ} methods or \acs{LINQ} query expressions,
which should be readable to any C\# programmer.
Although the fact that even parameters that do not filter the results go into \lstinline{Where()} can be confusing.

The syntax for \texttt{prop} modules, with \acs{LINQ} queries inside \lstinline{Select()},
usually following another \acs{LINQ} query, can be quite complicated,
but we believe it describes the meaning of the code quite well,
so it should still be fairly readable.

\item The various actions possible through LinqToWiki are highly discoverable:
every action begins in a single point – the \lstinline{Wiki} class.
An \ac{IDE}, such as Microsoft Visual Studio,
can then show the user the available actions, along with their description,
through autocompletion (called IntelliSense in Visual Studio).

\item The whole library is strongly-typed:
accessing modules, setting their parameters and then accessing their results
never involves using string constants,
that would represent some module, parameter or result property.
This is because actual types with methods, method parameters and properties are used.

Thanks to this, the chance for user error when using this library is greatly lowered.

\item The library is flexible: If some module in the \ac{API} of a wiki changes
(which happens regularly), the library should be still usable.
This is achieved by regenerating code for the wiki.

The same principle also applies to different wikis.
If the modules in several wikis differ, code can be generated for each wiki separately.
\end{itemize}

\medskip

Probably the biggest difference between the original goal and the final library is naming.
Names of types, methods, method parameters and properties should follow the .NET naming guidelines,
but this was not done, because we could not separate names from the \ac{API} into words.

\section*{Sample queries}

Following are samples of queries of different modules using LinqToWiki.
They show how some of the goals of this work have been achieved.

\begin{lstlisting}
string diff = wiki.compare(fromrev: 486474789,
                              torev: 487063697)
                    .value;
\end{lstlisting}

This query compares the text of two revisions and returns differences between them.
It shows that queries for modules that return a single item take their parameters as named method parameters
and directly return their result.

\begin{lstlisting}
var pages = wiki.Query.querypage(querypagepage.Uncategorizedpages)
                  .ToList();
\end{lstlisting}

This query returns a list of uncategorized pages,
as provided by the special page \texttt{Special:UncategorizedPages}.
It shows that required parameters to list-returning modules are given as method parameters
and that \lstinline{ToList()} (or, alternatively, \lstinline{ToEnumerable()})
has to be called on the result before it can be used.

\begin{lstlisting}
var pages = (from cm in wiki.Query.categorymembers()
	       where cm.title == "Category:Query languages"
	           && cm.type == categorymemberstype.subcat
	       select cm).Pages;
\end{lstlisting}

This query returns subcategories of the category \texttt{Query languages}.
The result is not directly usable, but it can be used in \texttt{prop} queries.
The query shows how LINQ can be used for querying list-returning modules
and that the \lstinline{Pages} property is used to create source for further queries.

\begin{lstlisting}
var result = pages.Select(
	p =>
	new
	{
		Info = p.info,
		Images = p.images().ToEnumerable()
	})
	.ToEnumerable();
\end{lstlisting}

This query returns information about images present on a list of pages specified by \lstinline{pages}
(see previous query).
It shows how the result of one query can be used as a source for a \texttt{prop} query,
how the LINQ method \lstinline{Select()} is used in such queries,
how anonymous types can be used to return the required information
and how \lstinline{ToEnumerable()} (or \lstinline{ToList()}) has to be also used inside \texttt{prop} queries.

%%% Seznam použité literatury
%%% Seznam použité literatury je zpracován podle platných standardů. Povinnou citační
%%% normou pro bakalářskou práci je ISO 690. Jména časopisů lze uvádět zkrácená, ale jen
%%% v kodifikované podobě. Všechny použité zdroje a prameny musí být řádně citovány.

\def\bibname{Bibliography}
\begin{thebibliography}{12}
\addcontentsline{toc}{chapter}{\bibname}

%\bibitem{lamport94}
%  {\sc Lamport,} Leslie.
%  \emph{\LaTeX: A Document Preparation System}.
%  2. vydání.
%  Massachusetts: Addison Wesley, 1994.
%  ISBN 0-201-52983-1.

\bibitem{mediawiki}
 \emph{MediaWiki}.
 MediaWiki.org.
 \url{http://www.mediawiki.org/wiki/MediaWiki}.

\bibitem{mediawiki-api}
 \emph{MediaWiki API}.
 MediaWiki.org.
 \url{http://www.mediawiki.org/wiki/API}.
 
\bibitem{cs-in-depth}
  {\sc Skeet,} Jon.
  \emph{C\# in Depth}.
  2nd edition.
  Stamford: Manning, 2011.
  Part 3, C\#~3: Revolutionizing how we code.
  ISBN 978-1-935182-47-4.
  
\bibitem{roslyn}
 \emph{Microsoft “Roslyn” CTP}.
 MSDN.
 \url{http://msdn.microsoft.com/en-US/roslyn}.
 
\bibitem{linq-to-sql-functions}
 \emph{Data Types and Functions}.
 LINQ to SQL Reference, MSDN Library.
 \url{http://msdn.microsoft.com/en-us/library/bb386970}.

\bibitem{warren}
 {\sc Warren}, Matt.
 \emph{LINQ: Building an IQueryable Provider -- Part III}.
 The Wayward WebLog.
 \url{http://blogs.msdn.com/b/mattwar/archive/2007/08/01/linq-building-an-iqueryable-provider-part-iii.aspx},
 2 August 2007.
 
\bibitem{reflection-emit}
 \emph{Emitting Dynamic Methods and Assemblies}.
 MSDN Library.
 \url{http://msdn.microsoft.com/en-us/library/8ffc3x75}.

\bibitem{codedom}
 \emph{Dynamic Source Code Generation and Compilation}.
 MSDN Library.
 \url{http://msdn.microsoft.com/en-us/library/650ax5cx}.
 
\bibitem{guidelines-for-names}
 \emph{Guidelines for Names}.
 Design Guidelines for Developing Class Libraries, MSDN Library.
 \url{http://msdn.microsoft.com/en-us/library/ms229002}.
 
\bibitem{wikitools}
 \emph{python-wikitools}.
 Google Code.
 \url{http://code.google.com/p/python-wikitools/}.
 
\bibitem{wikifunctions}
 \emph{WikiFunctions}.
 Wikipedia.
 \url{http://en.wikipedia.org/wiki/Wikipedia:WikiFunctions}.
 
\bibitem{linq-to-wikipedia}
 \emph{Linq to Wikipedia}.
 CodePlex.
 \url{http://linqtowikipedia.codeplex.com/}.
 
\end{thebibliography}

%%% Tabulky v bakalářské práci, existují-li.
\listoffigures
\addcontentsline{toc}{chapter}{List of Figures}

%%% Použité zkratky v bakalářské práci, existují-li, včetně jejich vysvětlení.
\chapwithtoc{List of Abbreviations}
\begin{acronym}[CodeDOM]
\acro{API}{Application Programming Interface}
\acro{CodeDOM}{Code Document Object Model}
\acro{CIL}{Common Intermediate Language}
\acro{CTP}{Community Technology Preview}
\acro{DLL}{Dynamic-Link Library}
\acro{DOM}{Document Object Model}
\acro{GUI}{Graphical User Interface}
\acro{HTTP}{Hypertext Transfer Protocol}
\acro{IDE}{Integrated Development Environment}
\acro{IP}{Internet Protocol}
\acro{JSON}{JavaScript Object Notation}
\acro{LINQ}{Language Integrated Query}
\acro{PHP}{PHP: Hypertext Preprocessor}
\acro{RSS}{Really Simple Syndication}
\acro{Rx}{Reactive Extensions}
\acro{SQL}{Structured Query Language}
\acro{TPL}{Task Parallel Library}
\acro{URL}{Uniform Resource Locator}
\acro{VB.NET}{Visual Basic .NET}
\acro{XML}{Extensible Markup Language}
\end{acronym}

%%% Přílohy k bakalářské práci, existují-li (různé dodatky jako výpisy programů,
%%% diagramy apod.). Každá příloha musí být alespoň jednou odkazována z vlastního
%%% textu práce. Přílohy se číslují.
\appendix
\chapter{Contents of the enclosed CD}
\label{cd}

The main part of the enclosed CD is the directory \texttt{LinqToWiki},
which contains all of the source code of the LinqToWiki library.
The same code is also available from the git repository of the library:
\url{https://github.com/svick/LINQ-to-Wiki/}.

The CD also contains this document (\texttt{thesis.pdf})
and a short presentation about the library (\texttt{presentation.pdf}).

\sectionnobreak{Using the library}

As described in Chapter~\ref{ltw}, using LinqToWiki consists of several steps.
Here, they are described in detail:

\begin{enumerate}
\item Compile the \texttt{lintowiki-codegen} application.

To do this, you need Microsoft Visual Studio 2012 with the September 2012 \ac{CTP} of Roslyn installed.
In it, open \texttt{LinqToWiki.sln} and build the project LinqToWiki.Codegen.App.

This step can be skipped, compiled version of \texttt{lintowiki-codegen} is included
in the \texttt{Tools} directory on the enclosed CD.

\item Run \texttt{lintowiki-codegen} to generate wiki-specific \ac{DLL}.

Execute the application \texttt{lintowiki-codegen.exe} with a parameter specifying which wiki
to use and optionally also other parameters altering the output.
The application requires .Net 4.5.
For further information, see Section~\ref{ltw-ca}.

This step can also be skipped, the enclosed CD contains \ac{DLL} generated for the English Wikipedia
in the directory \texttt{Lib\string\LinqToWiki.Generated}.

\item Use the generated \ac{DLL} in your application.

The LinqToWiki.Samples application (more details in Section~\ref{ltw-s}) is an example of application
that uses the generated \ac{DLL} to perform queries.

To build, execute and possibly modify it, open the solution file \texttt{LinqToWiki no codegen.sln}
in Microsoft Visual Studio 2010 or 2012.
\end{enumerate}

\openright
\end{document}